\documentclass{article}
\usepackage[margin=0.8in]{geometry}
\usepackage{amsmath,amssymb,amsthm,mathtools}
\usepackage{mathrsfs,bm,bbm}
\usepackage{graphicx,booktabs,array,multirow,tabularx,longtable}
\usepackage{enumitem}
\usepackage{xcolor}
\usepackage{algorithm,algpseudocode}
\usepackage{subcaption,tikz}
\usepackage{siunitx,cancel,accents}
\usepackage{microtype}
\usepackage[numbers,sort&compress]{natbib}
\usepackage[colorlinks=true,linkcolor=blue,citecolor=blue,urlcolor=blue]{hyperref}
\usepackage{cleveref}
\setlength{\emergencystretch}{2em}
\newtheorem{assumption}{Assumption}
\newtheorem{definition}{Definition}
\newtheorem{theorem}{Theorem}
\newtheorem{lemma}{Lemma}
\newtheorem{proposition}{Proposition}
\newtheorem{openendedquestion}{Open-ended Question}
\newtheorem{conjecture}{Conjecture}
\newcommand{\coreref}[1]{\ref{#1}}

\begin{document}

\sloppy

\section*{stat\_\allowbreak{}discrete\_\allowbreak{}optimal\_\allowbreak{}value\_\allowbreak{}minimax}

\noindent\textit{Specialization:} diagonal\_\allowbreak{}adaptive\_\allowbreak{}ratio\_\allowbreak{}tv\par

\paragraph{TL;DR.} For the identified scalar oracle optimal-policy value, the proved global finite-sample minimax bracket is \(c_\epsilon\min\{1,\sqrt d/n+d/[n\log(en)]\}\leq\mathfrak R_{n,d,\epsilon}\leq\min\{1,C_\epsilon d/n\}\). The upper bound is attained by an explicit one-pass empirical cellwise ratio estimator. An exact equal-propensity \(L_1\) Markov reduction, a paired-sign fuzzy mixture, and null-cell padding establish the converse. The two lower mechanisms cross at \(d\asymp\log^2(en)\), but this is not a matched minimax phase boundary. The same witness rules out the proposed direct four-cell lift of the JiaoHanWeissman2018L1 split-sample diagonal-adaptive certificate at \(n^{-1}+d/[n\log(en)]\), without disputing that source result in its normalized multinomial experiment; parametric MSE holds exactly when \(d=O(1)\), consistency requires \(d=o(n\log n)\), and \(d=o(n)\) is sufficient.

\section{Project justification}

\paragraph{Gap.} The accepted in-repository discrete-ATE theorem handles a linear target and obtains a matched \(n^{-1}+d^2/[n^2\log^2 n]\) MSE by a heavy/light hybrid. Local optimal-value regularity and gap-density smoothing results do not determine the unrestricted growing-alphabet point risk. Manski2004 and Stoye2009 study finite-sample minimax welfare regret of treatment allocations, not squared-error reporting of an observational scalar value. The nonlinear maximum adds a tie singularity, and the proposed direct four-cell extension of the JiaoHanWeissman2018L1 diagonal-adaptive workflow with a logarithmic certificate is false.

\paragraph{Niche.} The oracle value decomposes into the smooth homogeneous ratio sum \(\Lambda\) and the nonsmooth homogeneous absolute contrast \(\Gamma\). The exact equal-propensity slice turns the target into an affine transform of multinomial \(L_1\) distance, while a separate paired-sign mixture detects a small-alphabet obstruction that is invisible in the linear ATE problem.

\paragraph{Fill.} The note proves an honest finite-sample minimax bracket under the unchanged unrestricted class: an empirical ratio estimator gives \(C_\epsilon d/n\), while exact experiment reduction, the normalized-distance lower theorem, a direct fuzzy-mixture calculation, and alphabet padding give the global lower bound. An explicit \(d_n\asymp\log(en)\) witness rules out the proposed four-cell transplantation of JiaoHanWeissman2018L1's normalized-multinomial diagonal-adaptive certificate, without contradicting that source result in its own experiment. Matching achievability and the intermediate consistency region remain open.

\section{Related work}

The accepted in-repository result \texttt{stat\_discrete\_ate\_minimax\_loggap/thm:sharp-minimax-fixed-interior} proves a matched MSE \(n^{-1}+d^2/[n^2\log^2 n]\) for the linear ATE using a heavy/light ratio-polynomial hybrid. Its quantitative reciprocal-polynomial and factorial-moment calculations inform the smooth \(\Lambda\) branch, but they do not control the absolute-contrast \(\Gamma\) branch. The nonlinear oracle maximum has an unrestricted tie singularity, the additional \(\sqrt d/n\) paired-sign lower mechanism, the normalized-distance term \(d/[n\log(en)]\), and only \(d/n\) certified achievability; neither theorem implies the other. Manski, \cite{Manski2004}, studies finite-sample maximum welfare regret of covariate-conditional empirical-success allocation rules in randomized experiments, while Stoye, \cite{Stoye2009}, characterizes exact finite-sample minimax-regret treatment rules, including with finite covariates. Their loss is welfare regret of an allocated rule relative to the oracle allocation, possibly under rule randomization; it neither implies nor is implied by squared-error estimation of the scalar observational oracle value here. In particular, low regret need not report the scalar value, while a scalar estimate provides neither cellwise treatment assignments nor a regret bound. Zeng, Balakrishnan, Han, and Kennedy, \cite{ZengBalakrishnanHanKennedy2024Discrete}, study a linear treatment-specific mean rather than this maximum. Luedtke and van der Laan, \cite{LuedtkeVanderLaan2016OptimalValue}, Theorem 1, identify the local pathwise-differentiability boundary at treatment-effect ties but do not give a growing-alphabet finite-sample minimax bracket. Whitehouse, Chen, Austern, and Syrgkanis, \cite{WhitehouseChenAusternSyrgkanis2025}, obtain root-\(n\) smoothed inference under gap-density and nuisance-rate conditions that exclude the unrestricted tie and occupancy configurations used here. Jiao, Han, and Weissman, \cite{JiaoHanWeissman2018L1}, Construction 2 and Theorem 6, give a split-sample diagonal-adaptive estimator with \(d/[n\log n]\)-scale MSE in the normalized two-multinomial experiment. The negative theorem rules out the proposed direct lift of that workflow with the same certificate to the larger observational four-cell ratio model; it does not contradict their normalized-multinomial result. Their Theorem 3 and Lemma 33 supply the known-\(\mathbf Q\) lower theorem and fixed-sample comparison used after the exact equal-propensity embedding. Polynomial work by Wu and Yang, \cite{WuYang2016Entropy}, Jiao, Venkat, Han, and Weissman, \cite{JiaoVenkatHanWeissman2015Functionals}, and Cai and Low, \cite{CaiLow2011L1}, motivates approximation methods but does not supply the missing four-cell ratio-cone certificate. The new negative result is the explicit \(d_n\asymp\log(en)\) causal witness showing a \(\sqrt{\log(en)}\) factor over the proposed transplanted certificate.

\section{Setup}

Target (estimand / structural parameter): \(V^\star(P)=E_P[\max_{a\in\{0,1\}}E_P\{Y(a)\mid X\}]\).
Primary functional / estimator / diagnostic: \(V^\star(P)=\sum_{x=1}^dp_x\max\{\mu_{0x},\mu_{1x}\}=\frac12\sum_{x=1}^d\{\Lambda(q_x)+\Gamma(q_x)\}\), with the global finite-sample bracket \(c_\epsilon\min\{1,\sqrt d/n+d/[n\log(en)]\}\leq\mathfrak R_{n,d,\epsilon}\leq\min\{1,C_\epsilon d/n\}\) for every \(n\geq1\) and \(d\geq2\); in the nonsaturated range the lower terms balance at \(d\asymp\log^2(en)\).
\begin{itemize}
  \item \(n\) --- \texttt{sample size}; \(\mathbb N\); \texttt{index}
  \par\smallskip\noindent\emph{Definition.} \texttt{Number of observed units.}
  \item \(d\) --- \texttt{alphabet size}; \(\mathbb N\); \texttt{index}
  \par\smallskip\noindent\emph{Definition.} \texttt{Number of discrete confounder cells.}
  \item \(\epsilon\) --- \texttt{overlap constant}; \((0,1/2)\); \texttt{regime parameter}
  \par\smallskip\noindent\emph{Definition.} \texttt{Fixed lower bound on both treatment probabilities.}
  \item \(\rho_\epsilon\) --- \texttt{range constant}; \((0,\infty)\); \texttt{regime parameter}
  \par\smallskip\noindent\emph{Definition.} Positive constant depending only on \(\epsilon\) that bounds the proved range \(d\le \rho_\epsilon n\log(en)\).
  \item \(c_\epsilon\) --- \texttt{minimax lower constant}; \((0,\infty)\); \texttt{theorem constant}
  \par\smallskip\noindent\emph{Definition.} Positive lower-bound constant depending only on \(\epsilon\) whose existence is asserted by the risk statements.
  \item \(C_\epsilon\) --- \texttt{minimax upper constant}; \((0,\infty)\); \texttt{theorem constant}
  \par\smallskip\noindent\emph{Definition.} Finite upper-bound constant depending only on \(\epsilon\) whose existence is asserted by the risk statements.
  \item \([d]\) --- \texttt{covariate alphabet}; \texttt{finite sets}; \texttt{index set}
  \par\smallskip\noindent\emph{Definition.} \([d]=\{1,\ldots,d\}\).
  \item \(a\) --- \texttt{treatment index}; \(\{0,1\}\); \texttt{index}
  \par\smallskip\noindent\emph{Definition.} \texttt{Generic treatment level.}
  \item \(y\) --- \texttt{outcome index}; \(\{0,1\}\); \texttt{index}
  \par\smallskip\noindent\emph{Definition.} \texttt{Generic outcome level.}
  \item \(x\) --- \texttt{cell index}; \([d]\); \texttt{index}
  \par\smallskip\noindent\emph{Definition.} \texttt{Generic confounder cell.}
  \item \(i\) --- \texttt{unit index}; \(\{1,\ldots,n\}\); \texttt{index}
  \par\smallskip\noindent\emph{Definition.} \texttt{Generic observed-\allowbreak{}unit index.}
  \item \(X\) --- \texttt{covariate}; \([d]\); \texttt{data}
  \par\smallskip\noindent\emph{Definition.} \texttt{Observed discrete confounder.}
  \item \(A\) --- \texttt{treatment}; \(\{0,1\}\); \texttt{data}
  \par\smallskip\noindent\emph{Definition.} \texttt{Observed binary treatment.}
  \item \(Y\) --- \texttt{outcome}; \(\{0,1\}\); \texttt{data}
  \par\smallskip\noindent\emph{Definition.} \texttt{Observed binary outcome.}
  \item \(Y(a)\) --- \texttt{potential outcome}; \(\{0,1\}\); \texttt{causal object}
  \par\smallskip\noindent\emph{Definition.} Outcome under treatment level \(a\in\{0,1\}\).
  \item \(O_i\) --- \texttt{observed unit}; \([d]\times\{0,1\}^2\); \texttt{sample}
  \par\smallskip\noindent\emph{Definition.} \(O_i=(X_i,A_i,Y_i)\) for \(i\in\{1,\ldots,n\}\).
  \item \(P\) --- \texttt{full-\allowbreak{}data law}; laws of \((X,A,Y,Y(0),Y(1))\); \texttt{model}
  \par\smallskip\noindent\emph{Definition.} A law whose observed margin is the distribution of \(O=(X,A,Y)\).
  \item \(p_x\) --- \texttt{cell mass}; \([0,1]\); \texttt{primitive}
  \par\smallskip\noindent\emph{Definition.} \(p_x=P(X=x)\).
  \item \(\pi_x\) --- \texttt{propensity score}; \([0,1]\); \texttt{primitive}
  \par\smallskip\noindent\emph{Definition.} \(\pi_x=P(A=1\mid X=x)\) when \(p_x>0\).
  \item \(\mu_{ax}\) --- \texttt{outcome regression}; \([0,1]\); \texttt{primitive}
  \par\smallskip\noindent\emph{Definition.} \(\mu_{ax}=E_P[Y\mid A=a,X=x]\) when \(p_x>0\).
  \item \(q_{ay,x}\) --- \texttt{observed joint cell mass}; \([0,1]\); \texttt{primitive}
  \par\smallskip\noindent\emph{Definition.} \(q_{ay,x}=P(X=x,A=a,Y=y)\).
  \item \(\psi_1(P)\) --- \texttt{treatment-\allowbreak{}specific mean}; \([0,1]\); \texttt{published comparator target}
  \par\smallskip\noindent\emph{Definition.} \(\psi_1(P)=\sum_{x=1}^dp_x\mu_{1x}\).
  \item \(a^\star(x)\) --- \texttt{identified oracle treatment rule}; maps from \([d]\) to \(\{0,1\}\); \texttt{causal decision}
  \par\smallskip\noindent\emph{Definition.} \(a^\star(x)=\min\arg\max_{a\in\{0,1\}}\mu_{ax}\) when \(p_x>0\), and \(a^\star(x)=0\) when \(p_x=0\).
  \item \(q_x\) --- \texttt{four-\allowbreak{}cell mass vector}; \(\mathbb R_+^4\); \texttt{primitive}
  \par\smallskip\noindent\emph{Definition.} \(q_x=(q_{00,x},q_{01,x},q_{10,x},q_{11,x})\), ordered by \((A,Y)=(0,0),(0,1),(1,0),(1,1)\).
  \item \(\Delta_x\) --- \texttt{cell treatment contrast}; \([-1,1]\); \texttt{primitive}
  \par\smallskip\noindent\emph{Definition.} \(\Delta_x=\mu_{1x}-\mu_{0x}\).
  \item \(\Delta_d\) --- \texttt{probability simplex}; subsets of \(\mathbb R_+^d\); \texttt{comparator parameter space}
  \par\smallskip\noindent\emph{Definition.} \(\Delta_d=\{R\in\mathbb R_+^d:\sum_{x=1}^dR_x=1\}\).
  \item \(\mathbf P\) --- \texttt{first multinomial distribution}; \(\Delta_d\); \texttt{lower-\allowbreak{}bound input}
  \par\smallskip\noindent\emph{Definition.} \(\mathbf P=(P_1,\ldots,P_d)\) is the first distribution in the normalized \(L_1\)-distance comparator experiment.
  \item \(\mathbf Q\) --- \texttt{second multinomial distribution}; \(\Delta_d\); \texttt{lower-\allowbreak{}bound input}
  \par\smallskip\noindent\emph{Definition.} \(\mathbf Q=(Q_1,\ldots,Q_d)\) is the second distribution in the normalized \(L_1\)-distance comparator experiment.
  \item \(\mathcal C_\epsilon\) --- \texttt{overlap cone}; subsets of \(\mathbb R_+^4\); \texttt{approximation domain}
  \par\smallskip\noindent\emph{Definition.} Cone constructed in \(\mathrm{def:overlap-cone}\).
  \item \(\Lambda\) --- \texttt{degree-\allowbreak{}one homogeneous ratio-\allowbreak{}sum functional}; \(\mathcal C_\epsilon\to[0,\infty)\); \texttt{functional}
  \par\smallskip\noindent\emph{Definition.} Continuous cell functional, smooth on \(\mathcal C_\epsilon\setminus\{0\}\), constructed in \(\mathrm{def:cell-functionals}\).
  \item \(\Gamma\) --- \texttt{homogeneous ratio-\allowbreak{}absolute-\allowbreak{}contrast}; \(\mathcal C_\epsilon\to[0,\infty)\); \texttt{functional}
  \par\smallskip\noindent\emph{Definition.} Continuous degree-one homogeneous cell functional constructed in \(\mathrm{def:cell-functionals}\).
  \item \(V^\star(P)\) --- \texttt{oracle optimal-\allowbreak{}policy value}; \([0,1]\); \texttt{target}
  \par\smallskip\noindent\emph{Definition.} Causal target identified in \(\mathrm{def:oracle-value}\).
  \item \(\mathcal P_{d,\epsilon}\) --- \texttt{discrete-\allowbreak{}confounder law class}; \texttt{law classes}; \texttt{model class}
  \par\smallskip\noindent\emph{Definition.} Class constructed in \(\mathrm{def:model-class}\).
  \item \(I_0\) --- \texttt{pilot sample block}; subsets of \(\{1,\ldots,n\}\); \texttt{construction input}
  \par\smallskip\noindent\emph{Definition.} \texttt{First block in a deterministic three-\allowbreak{}way partition with block sizes differing by at most one.}
  \item \(I_1\) --- \texttt{ratio-\allowbreak{}estimation sample block}; subsets of \(\{1,\ldots,n\}\); \texttt{construction input}
  \par\smallskip\noindent\emph{Definition.} \texttt{Second block in the deterministic three-\allowbreak{}way partition.}
  \item \(I_2\) --- \texttt{factorial-\allowbreak{}estimation sample block}; subsets of \(\{1,\ldots,n\}\); \texttt{construction input}
  \par\smallskip\noindent\emph{Definition.} \texttt{Third block in the deterministic three-\allowbreak{}way partition.}
  \item \(N^{(j)}_{ay,x}\) --- \texttt{split joint cell count}; \(\mathbb N_0\); \texttt{statistic}
  \par\smallskip\noindent\emph{Definition.} \(N^{(j)}_{ay,x}=\sum_{i\in I_j}\mathbf 1\{X_i=x,A_i=a,Y_i=y\}\) for \(j\in\{0,1,2\}\).
  \item \(r\) --- \texttt{four-\allowbreak{}dimensional multi-\allowbreak{}index}; \(\mathbb N_0^4\); \texttt{construction input}
  \par\smallskip\noindent\emph{Definition.} Exponent vector \(r=(r_{00},r_{01},r_{10},r_{11})\) with \(\lvert r\rvert=\sum_{a,y}r_{ay}\).
  \item \(K_n\) --- \texttt{polynomial degree}; \(\mathbb N\); \texttt{tuning parameter}
  \par\smallskip\noindent\emph{Definition.} \(K_n=\max\{1,\lfloor\epsilon^2\log(en)/2^{16}\rfloor\}\).
  \item \(\widehat{\mathcal R}_n\) --- \texttt{pilot regular region}; subsets of \([d]\); \texttt{statistic}
  \par\smallskip\noindent\emph{Definition.} Pilot-certified massive cells separated from the tie diagonal, constructed in \(\mathrm{def:pilot-partition}\).
  \item \(\widehat{\mathcal D}_n\) --- \texttt{pilot diagonal-\allowbreak{}or-\allowbreak{}sparse region}; subsets of \([d]\); \texttt{statistic}
  \par\smallskip\noindent\emph{Definition.} Complement of the pilot regular region, constructed in \(\mathrm{def:pilot-partition}\).
  \item \(\widetilde B_{x,n}\) --- \texttt{raw pilot-\allowbreak{}local feasible region}; compact subsets of \(\mathcal C_\epsilon\); \texttt{approximation domain}
  \par\smallskip\noindent\emph{Definition.} Possibly empty confidence-box intersection constructed in \(\mathrm{def:pilot-partition}\).
  \item \(B_{x,n}\) --- \texttt{total pilot-\allowbreak{}local overlap region}; nonempty compact subsets of \(\mathcal C_\epsilon\); \texttt{approximation domain}
  \par\smallskip\noindent\emph{Definition.} Raw confidence-box intersection with the deterministic singleton fallback constructed in \(\mathrm{def:pilot-partition}\).
  \item \(\widehat V_n^{\mathrm{DAT}}\) --- \texttt{empirical cellwise ratio fallback estimator}; \([0,1]\); \texttt{estimator}
  \par\smallskip\noindent\emph{Definition.} Full-sample plug-in statistic constructed in \(\mathrm{def:diagonal-adaptive-estimator}\).
  \item \(\mathfrak R_{n,d,\epsilon}\) --- \texttt{minimax mean-\allowbreak{}squared error}; \([0,1]\); \texttt{risk}
  \par\smallskip\noindent\emph{Definition.} Risk constructed in \(\mathrm{def:minimax-risk}\).
  \item \(\mathsf K_{n,d}\) --- \texttt{Markov experiment-\allowbreak{}reduction kernel}; \texttt{Markov kernels from two multinomial samples to causal samples}; \texttt{lower-\allowbreak{}bound construction}
  \par\smallskip\noindent\emph{Definition.} Kernel constructed in \(\mathrm{def:l1-embedding}\) from \(n\) draws from each of \(\mathbf P\) and \(\mathbf Q\) to \(n\) observed causal triples.
  \item \(\mathfrak R^{L_1}_{n,d}\) --- \texttt{normalized multinomial distance minimax mean-\allowbreak{}squared error}; \([0,4]\); \texttt{comparator risk}
  \par\smallskip\noindent\emph{Definition.} Risk constructed in \(\mathrm{def:l1-embedding}\) for estimating \(\lVert\mathbf P-\mathbf Q\rVert_1\) from \(n\) samples from each distribution.
  \item \(\mathscr H_\epsilon^{\mathrm{dual}}\) --- \texttt{exact-\allowbreak{}constant derivation handle}; \texttt{approximation-\allowbreak{}duality programs}; \texttt{open construction}
  \par\smallskip\noindent\emph{Definition.} Handle constructed in \(\mathrm{def:exact-constant-handle}\).
\end{itemize}

\section{Assumptions}

\begin{assumption}[ass:iid-sampling]\label{ass:iid-sampling}
\(O_1,\ldots,O_n\) are independent draws from the observed \((X,A,Y)\)-margin of \(P\).
\par\smallskip\noindent\emph{Standard} (independent and identically distributed sampling, \cite{ZengBalakrishnanHanKennedy2024Discrete}).
\end{assumption}

\begin{assumption}[ass:consistency]\label{ass:consistency}
\(Y=Y(A)\) almost surely.
\par\smallskip\noindent\emph{Standard} (consistency, \cite{ZengBalakrishnanHanKennedy2024Discrete}).
\end{assumption}

\begin{assumption}[ass:conditional-exchangeability]\label{ass:conditional-exchangeability}
\((Y(0),Y(1))\perp A\mid X\).
\par\smallskip\noindent\emph{Standard} (conditional exchangeability, \cite{ZengBalakrishnanHanKennedy2024Discrete}).
\end{assumption}

\begin{assumption}[ass:fixed-overlap]\label{ass:fixed-overlap}
\(\epsilon\le\pi_x\le1-\epsilon\) for every \(x\in[d]\) with \(p_x>0\).
\par\smallskip\noindent\emph{Standard} (fixed strong overlap, \cite{ZengBalakrishnanHanKennedy2024Discrete}).
\end{assumption}

\section{Definitions}

\begin{definition}[def:model-class]\label{def:model-class}
\par\smallskip\noindent\emph{Defined object.} \texttt{observationally equivalent full-\allowbreak{}data completion of the published unrestricted discrete-\allowbreak{}confounder observed-\allowbreak{}data experiment}
\par\smallskip\noindent\emph{Construction.} \(\mathcal P_{d,\epsilon}:=\{P: P\text{ is a law of }(X,A,Y,Y(0),Y(1))\text{ with }X\in[d],\ A,Y,Y(0),Y(1)\in\{0,1\},\text{ and }P\text{ satisfies }\mathrm{ass:consistency},\mathrm{ass:conditional-exchangeability},\mathrm{ass:fixed-overlap}\}\) for every \(d\in\mathbb N\) and \(0<\epsilon<1/2\).
\par\smallskip\noindent\emph{Member properties.} \texttt{ass:\allowbreak{}consistency}; \texttt{ass:\allowbreak{}conditional-\allowbreak{}exchangeability}; \texttt{ass:\allowbreak{}fixed-\allowbreak{}overlap}.
\end{definition}

\begin{definition}[def:overlap-cone]\label{def:overlap-cone}
\par\smallskip\noindent\emph{Defined object.} \texttt{overlap-\allowbreak{}restricted four-\allowbreak{}cell cone}
\par\smallskip\noindent\emph{Construction.} For \(u=(u_{00},u_{01},u_{10},u_{11})\in\mathbb R_+^4\), put \(s_a(u)=u_{a0}+u_{a1}\) and \(s(u)=s_0(u)+s_1(u)\). Define \(\mathcal C_\epsilon:=\{u\in\mathbb R_+^4:\epsilon s(u)\le s_1(u)\le(1-\epsilon)s(u)\}\).
\par\smallskip\noindent\emph{Inputs.} \(\epsilon\).
\end{definition}

\begin{definition}[def:cell-functionals]\label{def:cell-functionals}
\par\smallskip\noindent\emph{Defined object.} \texttt{homogeneous ratio sum and ratio-\allowbreak{}absolute-\allowbreak{}contrast}
\par\smallskip\noindent\emph{Construction.} For nonzero \(u\in\mathcal C_\epsilon\), define \(\Lambda(u):=s(u)\{u_{11}/s_1(u)+u_{01}/s_0(u)\}\) and \(\Gamma(u):=s(u)\lvert u_{11}/s_1(u)-u_{01}/s_0(u)\rvert\); set \(\Lambda(0)=\Gamma(0)=0\). For every \(t\ge0\), \(\Lambda(tu)=t\Lambda(u)\) and \(\Gamma(tu)=t\Gamma(u)\), with \(0\le\Lambda(u)\le2s(u)\) and \(0\le\Gamma(u)\le s(u)\). Thus both maps have linear growth on the unrestricted cone rather than a bounded range. The map \(\Lambda\) is smooth only on \(\mathcal C_\epsilon\setminus\{0\}\); there \(\Gamma\) is nonsmooth exactly on the ratio-tie diagonal \(u_{11}s_0(u)=u_{01}s_1(u)\).
\par\smallskip\noindent\emph{Inputs.} \(\mathcal C_\epsilon\).
\end{definition}

\begin{definition}[def:oracle-value]\label{def:oracle-value}
\par\smallskip\noindent\emph{Defined object.} \texttt{identified unrestricted oracle value}
\par\smallskip\noindent\emph{Construction.} \(V^\star(P):=E_P[\max_{a\in\{0,1\}}E_P\{Y(a)\mid X\}]=\sum_{x=1}^d p_x\max\{\mu_{0x},\mu_{1x}\}=\frac12\sum_{x=1}^d\{\Lambda(q_x)+\Gamma(q_x)\}\).
\par\smallskip\noindent\emph{Inputs.} \texttt{ass:\allowbreak{}consistency}; \texttt{ass:\allowbreak{}conditional-\allowbreak{}exchangeability}; \texttt{ass:\allowbreak{}fixed-\allowbreak{}overlap}; \texttt{def:\allowbreak{}cell-\allowbreak{}functionals}.
\end{definition}

\begin{definition}[def:pilot-partition]\label{def:pilot-partition}
\par\smallskip\noindent\emph{Defined object.} \texttt{mass-\allowbreak{}and-\allowbreak{}tie pilot partition}
\par\smallskip\noindent\emph{Construction.} Let \(m_j=\lvert I_j\rvert\), set \(\widehat q^{(0)}_{ay,x}=N^{(0)}_{ay,x}/\max\{m_0,1\}\), and put \(\widehat p_x^{(0)}=\sum_{a,y}\widehat q^{(0)}_{ay,x}\). Let \(\widehat\Delta_x^{(0)}\) be the difference of the two empirical arm means, with \(0/0=0\). Define \(\widehat{\mathcal R}_n:=\{x:\widehat p_x^{(0)}\ge256\log(en)/\max\{m_0,1\},\ \min_a\sum_yN^{(0)}_{ay,x}\ge64\log(en),\ \lvert\widehat\Delta_x^{(0)}\rvert\ge32\sqrt{\log(en)/(\max\{m_0,1\}\widehat p_x^{(0)})}\}\) and \(\widehat{\mathcal D}_n=[d]\setminus\widehat{\mathcal R}_n\), interpreting the last inequality as false when \(\widehat p_x^{(0)}=0\). For every \(x\in[d]\), put \(w_{ay,x}=16\sqrt{\widehat q^{(0)}_{ay,x}\log(en)/\max\{m_0,1\}}+16\log(en)/\max\{m_0,1\}\) and \(\widetilde B_{x,n}:=\mathcal C_\epsilon\cap\prod_{a,y}[\max\{0,\widehat q^{(0)}_{ay,x}-w_{ay,x}\},\widehat q^{(0)}_{ay,x}+w_{ay,x}]\). Define the total compact-box convention \(B_{x,n}:=\widetilde B_{x,n}\) when \(\widetilde B_{x,n}\ne\varnothing\), and \(B_{x,n}:=\{0\}\) otherwise. This definition applies on every sample outcome and to cells in both pilot regions.
\par\smallskip\noindent\emph{Inputs.} \(O_i\); \(I_0\); \(I_1\); \(I_2\); \(N^{(j)}_{ay,x}\); \(\mathcal C_\epsilon\).
\end{definition}

\begin{definition}[def:diagonal-adaptive-estimator]\label{def:diagonal-adaptive-estimator}
\par\smallskip\noindent\emph{Defined object.} \texttt{explicit empirical cellwise ratio fallback estimator}
\par\smallskip\noindent\emph{Construction.} Given \(O_1,\ldots,O_n\), for every \(x\in[d]\) and \(a\in\{0,1\}\) define
\[
 \overline N_x=\sum_{i=1}^n\mathbf 1\{X_i=x\},\qquad
 \overline N_{ax}=\sum_{i=1}^n\mathbf 1\{X_i=x,A_i=a\},\qquad
 \overline S_{ax}=\sum_{i=1}^n\mathbf 1\{X_i=x,A_i=a,Y_i=1\}.
\]
Put \(\widehat p_x=\overline N_x/n\) and \(\widehat\mu_{ax}=\overline S_{ax}/\overline N_{ax}\), with \(0/0=0\).  Define the full-sample empirical cellwise ratio estimator
\[
 \widehat V_n^{\mathrm{DAT}}
 :=\sum_{x=1}^d\widehat p_x
       \max_{a\in\{0,1\}}\widehat\mu_{ax}.
\]
It is a finite, one-pass computable statistic and belongs to \([0,1]\).  The legacy definition identifier is retained only to preserve stable references; this construction replaces the uncertified pilot-polynomial selector.
\par\smallskip\noindent\emph{Inputs.} \(O_i\).
\end{definition}

\begin{definition}[def:minimax-risk]\label{def:minimax-risk}
\par\smallskip\noindent\emph{Defined object.} \texttt{minimax oracle-\allowbreak{}value risk}
\par\smallskip\noindent\emph{Construction.} \(\mathfrak R_{n,d,\epsilon}:=\inf_{\widehat V}\sup_{P\in\mathcal P_{d,\epsilon}}E_P[(\widehat V(O_1,\ldots,O_n)-V^\star(P))^2]\), where the infimum is over all measurable estimators.
\par\smallskip\noindent\emph{Inputs.} \texttt{def:\allowbreak{}model-\allowbreak{}class}; \texttt{def:\allowbreak{}oracle-\allowbreak{}value}.
\end{definition}

\begin{definition}[def:l1-embedding]\label{def:l1-embedding}
\par\smallskip\noindent\emph{Defined object.} exact normalized equal-propensity \(L_1\) embedding and Markov reduction
\par\smallskip\noindent\emph{Construction.} For every \((\mathbf P,\mathbf Q)\in\Delta_d\times\Delta_d\), define an observed causal law by \(p_x=(P_x+Q_x)/2\), \(\pi_x=1/2\), \(\mu_{1x}=P_x/(P_x+Q_x)\), and \(\mu_{0x}=Q_x/(P_x+Q_x)\), setting both ratios to zero when \(P_x+Q_x=0\). Equivalently, \(q_{11,x}=q_{00,x}=P_x/4\) and \(q_{10,x}=q_{01,x}=Q_x/4\). Complete this observed law to a member of \(\mathcal P_{d,\epsilon}\) by taking \(Y(0)\) and \(Y(1)\) conditionally independent Bernoulli variables with means \(\mu_{0x}\) and \(\mu_{1x}\), drawing \(A\mid X=x\sim\operatorname{Bernoulli}(1/2)\) independently of \((Y(0),Y(1))\) conditional on \(X\), and setting \(Y=Y(A)\). The Markov kernel \(\mathsf K_{n,d}\) takes independent streams \(U_1,\ldots,U_n\sim\mathbf P\) and \(W_1,\ldots,W_n\sim\mathbf Q\); independently for every \(i\), it tosses two fair coins \(S_i,D_i\), uses \(X_i=U_i\) when \(S_i=1\) and \(X_i=W_i\) when \(S_i=0\), assigns \((A_i,Y_i)=(1,1)\) or \((0,0)\) according to \(D_i\) on the first stream, and assigns \((A_i,Y_i)=(1,0)\) or \((0,1)\) according to \(D_i\) on the second stream. Finally, \(\mathfrak R^{L_1}_{n,d}:=\inf_{\widehat L}\sup_{(\mathbf P,\mathbf Q)\in\Delta_d^2}E[(\widehat L-\lVert\mathbf P-\mathbf Q\rVert_1)^2]\), where \(\widehat L\) uses the two input streams.
\par\smallskip\noindent\emph{Inputs.} \texttt{def:\allowbreak{}model-\allowbreak{}class}; \texttt{def:\allowbreak{}oracle-\allowbreak{}value}; \(\mathbf P\); \(\mathbf Q\).
\end{definition}

\begin{definition}[def:exact-constant-handle]\label{def:exact-constant-handle}
\par\smallskip\noindent\emph{Defined object.} \texttt{approximation-\allowbreak{}duality handle for the leading constant}
\par\smallskip\noindent\emph{Construction.} \(\mathscr H_\epsilon^{\mathrm{dual}}\) is the program that first rescales each pilot-local cone by its Poisson noise radius, then computes the limiting best-polynomial error for \(\Gamma\), and finally compares its moment-problem dual with the normalized equal-propensity \(L_1\) experiment in \(\mathrm{def:l1-embedding}\). This handle specifies the route to the local asymptotic minimax constant without fixing its value.
\par\smallskip\noindent\emph{Inputs.} \texttt{def:\allowbreak{}pilot-\allowbreak{}partition}; \texttt{def:\allowbreak{}cell-\allowbreak{}functionals}; \texttt{def:\allowbreak{}l1-\allowbreak{}embedding}.
\end{definition}

\section{Main results}

\begin{proposition}[prop:identification-decomposition]\label{prop:identification-decomposition}
Under \(\mathrm{ass:consistency}\), \(\mathrm{ass:conditional-exchangeability}\), and \(\mathrm{ass:fixed-overlap}\), the rule \(a^\star(x)=\min\arg\max_{a\in\{0,1\}}\mu_{ax}\) is identified on \(\{x:p_x>0\}\), is defined as zero when \(p_x=0\), and has value \(V^\star(P)=\frac12\sum_{x=1}^d\{\Lambda(q_x)+\Gamma(q_x)\}\).
\end{proposition}
\begin{proof}
\textbf{Proof.}
Fix \(P\in\mathcal P_{d,\epsilon}\) and a cell \(x\) with \(p_x>0\).  Fixed overlap gives
\[
 P(A=a\mid X=x)>0,
 \qquad a\in\{0,1\},
\]
so both observed conditional means \(\mu_{ax}\) are well defined.  By conditional exchangeability and consistency,
\[
 E_P\{Y(a)\mid X=x\}
 =E_P\{Y(a)\mid A=a,X=x\}
 =E_P\{Y\mid A=a,X=x\}
 =\mu_{ax}.
\]
Thus the two conditional causal means, their ordering, and the tie-broken maximizer
\(a^\star(x)=\min\arg\max_a\mu_{ax}\) are identified from the observed law on the support \(\{x:p_x>0\}\).  The stipulated value \(a^\star(x)=0\) on a null cell is immaterial to every value functional.

For the four-cell vector \(q_x\), write \(s_a(q_x)=q_{a0,x}+q_{a1,x}\) and \(s(q_x)=s_0(q_x)+s_1(q_x)\).  Directly from the definitions,
\[
 s(q_x)=p_x,
 \qquad s_1(q_x)=p_x\pi_x,
 \qquad s_0(q_x)=p_x(1-\pi_x).
\]
Fixed overlap makes the last two quantities positive.  Consequently,
\[
 \frac{q_{11,x}}{s_1(q_x)}=\mu_{1x},
 \qquad
 \frac{q_{01,x}}{s_0(q_x)}=\mu_{0x}.
\]
The definitions of \(\Lambda\) and \(\Gamma\) therefore give
\[
 \Lambda(q_x)=p_x(\mu_{1x}+\mu_{0x}),
 \qquad
 \Gamma(q_x)=p_x\lvert\mu_{1x}-\mu_{0x}\rvert.
\]
Using \(\max\{u,v\}=\{u+v+\lvert u-v\rvert\}/2\),
\[
 p_x\max\{\mu_{0x},\mu_{1x}\}
 =\frac12\{\Lambda(q_x)+\Gamma(q_x)\}.
\]
When \(p_x=0\), all four coordinates of \(q_x\) vanish, and both sides are zero by the conventions in \(\mathrm{def:cell-functionals}\).  Summing over \(x\in[d]\) and invoking \(\mathrm{def:oracle-value}\) proves
\[
 V^\star(P)=\frac12\sum_{x=1}^d\{\Lambda(q_x)+\Gamma(q_x)\}.
\]
\end{proof}
\par\smallskip\noindent \textit{Justification.} The decomposition isolates the ordinary ratio problem from the absolute-value cusp that changes the minimax frontier. \textit{Gap.} LuedtkeVanderLaan2016OptimalValue, Theorem 1, characterizes pathwise differentiability and the canonical gradient of the optimal value, but does not express the growing-alphabet experiment as a homogeneous four-cell approximation problem. \textit{Consumer.} The decomposition supplies the target used by the upper estimator and the normalized \(L_1\) lower-bound transfer.

\begin{lemma}[lem:l1-lower-transfer]\label{lem:l1-lower-transfer}
For every \((\mathbf P,\mathbf Q)\in\Delta_d^2\), the law produced by \(\mathrm{def:l1-embedding}\) belongs to \(\mathcal P_{d,\epsilon}\), satisfies \(\sum_{x=1}^dp_x=1\), and obeys \(V^\star(P)=1/2+\lVert\mathbf P-\mathbf Q\rVert_1/4\). The causal sample is the Markov image under \(\mathsf K_{n,d}\) of \(n\) independent \(\mathbf P\)-draws and \(n\) independent \(\mathbf Q\)-draws, so total variation contracts under this reduction. For every causal estimator \(\widehat V\), the composed estimator \(\widehat L=4\widehat V\circ\mathsf K_{n,d}-2\) therefore yields \(\mathfrak R_{n,d,\epsilon}\ge\mathfrak R^{L_1}_{n,d}/16\). A paired-sign fuzzy mixture gives \(\mathfrak R^{L_1}_{n,d}\ge c_\epsilon d/\{n\log(en)\}\) for \(2\le d\le C\log^2(en)\); for \(C\log^2(en)\le d\le\rho_\epsilon n\log(en)\), the same bound follows by specializing Theorem 3 of JiaoHanWeissman2018L1 to uniform \(\mathbf Q\) and checking its dominance condition. Independently, a regular two-point submodel gives \(\mathfrak R_{n,d,\epsilon}\ge c_\epsilon/n\).
\end{lemma}
\begin{proof}
\textbf{Proof.}
Fix \((\mathbf P,\mathbf Q)\in\Delta_d^2\).  From \(\mathrm{def:l1-embedding}\),
\[
 p_x=\sum_{a,y}q_{ay,x}=\frac{P_x+Q_x}{2},
 \qquad
 \sum_{x=1}^dp_x=1.
\]
On every occupied cell the treatment probability is \(1/2\), and the potential-outcome completion specified in that definition satisfies consistency and conditional exchangeability.  Since \(0<\epsilon<1/2\), fixed overlap holds.  Thus the completed law belongs to \(\mathcal P_{d,\epsilon}\).  By \(\mathrm{prop:identification-decomposition}\),
\[
 \begin{aligned}
 V^\star(P)
 &=\frac12\sum_{x=1}^d\max\{P_x,Q_x\}\\
 &=\frac14\sum_{x=1}^d
   \{P_x+Q_x+\lvert P_x-Q_x\rvert\}\\
 &=\frac12+\frac14\lVert\mathbf P-\mathbf Q\rVert_1.
 \end{aligned}
\]

For each observation, \(\mathsf K_{n,d}\) chooses one input stream with a fair coin and then one of that stream's two associated \((A,Y)\)-pairs with a fair coin.  Its output probabilities are exactly
\[
 q_{11,x}=q_{00,x}=\frac{P_x}{4},
 \qquad
 q_{10,x}=q_{01,x}=\frac{Q_x}{4}.
\]
Independence of the input draws and kernel coins across units therefore produces exactly \(n\) independent observations from the embedded causal law, as required by \(\mathrm{ass:iid-sampling}\).  For completeness, if \(K\) is any Markov kernel and \(\nu,\nu'\) are two input laws, then for every output event \(B\),
\[
 \lvert(\nu K)(B)-(\nu'K)(B)\rvert
 =\left\lvert\int K(z,B)\,d(\nu-\nu')(z)\right\rvert
 \le\lVert\nu-\nu'\rVert_{\mathrm{TV}},
\]
because \(0\le K(z,B)\le1\).  Taking the supremum over \(B\) proves total-variation contraction.

For a causal estimator \(\widehat V\), define the randomized two-stream estimator
\[
 \widehat L=4\widehat V\circ\mathsf K_{n,d}-2.
\]
The target identity gives, pointwise in \((\mathbf P,\mathbf Q)\),
\[
 E_{\mathbf P,\mathbf Q}
 \left[(\widehat L-\lVert\mathbf P-\mathbf Q\rVert_1)^2\right]
 =16E_P\left[(\widehat V-V^\star(P))^2\right].
\]
If the comparator infimum is restricted to nonrandomized estimators, conditional expectation given the two input streams derandomizes \(\widehat L\), and Jensen's inequality does not increase squared risk.  Taking suprema and infima proves
\[
 \mathfrak R_{n,d,\epsilon}\ge\frac1{16}\mathfrak R^{L_1}_{n,d}.
\]
These membership, value, and decision-theoretic conclusions also agree with \(\mathrm{lem:exact-l1-transfer-parametric-floor}\), which in addition gives, for every \(n\ge1\) and \(d\ge2\),
\[
 \mathfrak R_{n,d,\epsilon}\ge\frac1{1024n}.
\]

Apply \(\mathrm{lem:two-band-l1-lower}\) to the exact factor-sixteen transfer.  There are universal \(c,\rho>0\) such that
\[
 \mathfrak R_{n,d,\epsilon}
 \ge c\frac d{n\log(en)},
 \qquad
 2\le d\le\rho n\log(en).
\]
The proof of that helper is the paired-sign fuzzy-mixture argument when \(d\le C\log^2(en)\), and the uniform-\(\mathbf Q\) specialization of the cited Jiao--Han--Weissman theorem, including its fixed-sample comparison, above that threshold.  Renaming the positive constants as \(c_\epsilon\) and \(\rho_\epsilon\) is legitimate; the embedded propensity is exactly \(1/2\), so the only overlap boundary used is \(0<\epsilon<1/2\).  Together with the preceding two-point floor, this proves every assertion.
\end{proof}
\par\smallskip\noindent \textit{Justification.} The exact Markov kernel turns every causal estimator into a comparator estimator with the precise squared-loss factor sixteen. \textit{Gap.} The published comparator theorem does not itself establish causal class membership, the causal value identity, or the Markov reduction. \textit{Consumer.} The global theorem uses the transferred normalized-distance bound and parametric floor.

\begin{proposition}[prop:l1-reduction]\label{prop:l1-reduction}
For every \((\mathbf P,\mathbf Q)\in\Delta_d^2\), the exact embedding in \(\mathrm{def:l1-embedding}\) has \(p_x=(P_x+Q_x)/2\), \(\pi_x=1/2\), \(q_{11,x}=q_{00,x}=P_x/4\), and \(q_{10,x}=q_{01,x}=Q_x/4\), with \(\sum_{x=1}^dp_x=1\) and fixed overlap. On this normalized slice, \[V^\star(P)=\frac12+\frac14\lVert\mathbf P-\mathbf Q\rVert_1.\] The slice is a proper subexperiment of the full observational four-cell model, whose unknown arm masses and ratios are handled by \(\widehat V_n^{\mathrm{DAT}}\).
\end{proposition}
\begin{proof}
\textbf{Proof.}
Fix \((\mathbf P,\mathbf Q)\in\Delta_d^2\), and put \(s_x=P_x+Q_x\).  By \(\mathrm{def:l1-embedding}\),
\[
 p_x=\frac{s_x}{2},\qquad \pi_x=\frac12.
\]
If \(s_x>0\), the definitions of the observed joint cell masses and the embedded conditional means give
\[
\begin{aligned}
 q_{11,x}
 &=p_x\pi_x\mu_{1x}
 =\frac{s_x}{2}\frac12\frac{P_x}{s_x}
 =\frac{P_x}{4},
&
 q_{10,x}
 &=p_x\pi_x(1-\mu_{1x})
 =\frac{Q_x}{4},\\
 q_{01,x}
 &=p_x(1-\pi_x)\mu_{0x}
 =\frac{s_x}{2}\frac12\frac{Q_x}{s_x}
 =\frac{Q_x}{4},
&
 q_{00,x}
 &=p_x(1-\pi_x)(1-\mu_{0x})
 =\frac{P_x}{4}.
\end{aligned}
\]
If \(s_x=0\), nonnegativity implies \(P_x=Q_x=0\).  The zero-ratio convention in \(\mathrm{def:l1-embedding}\) then yields the same four identities.  Since \(\mathbf P\) and \(\mathbf Q\) both belong to \(\Delta_d\),
\[
 \sum_{x=1}^d p_x
 =\frac12\sum_{x=1}^d(P_x+Q_x)
 =1.
\]
Moreover, \(0<\epsilon<1/2\) implies
\[
 \epsilon\le \pi_x=\frac12\le 1-\epsilon
\]
on every supported cell, so \(\mathrm{ass:fixed-overlap}\), as incorporated in \(\mathrm{def:model-class}\), holds.

We next evaluate the causal target through its given identifying functional, rather than identifying the target by definition.  By \(\mathrm{def:oracle-value}\), cells with \(s_x=0\) contribute zero and
\[
\begin{aligned}
 V^\star(P)
 &=\sum_{x=1}^d \frac{P_x+Q_x}{2}
   \max\!\left\{\frac{Q_x}{P_x+Q_x},
                    \frac{P_x}{P_x+Q_x}\right\}\\
 &=\frac12\sum_{x=1}^d\max\{P_x,Q_x\}\\
 &=\frac14\sum_{x=1}^d
   \bigl(P_x+Q_x+\lvert P_x-Q_x\rvert\bigr)\\
 &=\frac12+\frac14\lVert\mathbf P-\mathbf Q\rVert_1.
\end{aligned}
\]
The third equality uses the algebraic identity \(\max\{u,v\}=(u+v+\lvert u-v\rvert)/2\), and the last equality again uses the simplex constraints.

It remains to verify properness of the slice.  For any \(d\ge1\), consider the observed law having \(p_1=1\), \(\pi_1=\epsilon\), and \(\mu_{01}=\mu_{11}=0\), with all other cells null.  Take both potential outcomes to be identically zero, draw \(A\mid X=1\) as \(\operatorname{Bernoulli}(\epsilon)\), and set \(Y=Y(A)\).  This completion satisfies consistency, conditional exchangeability, and fixed overlap, hence belongs to \(\mathcal P_{d,\epsilon}\).  Its supported-cell propensity is \(\epsilon\ne1/2\), whereas every embedded law has propensity \(1/2\) on every supported cell.  Thus the normalized slice is a strict subset of the full observational four-cell experiment.

Finally, \(\mathrm{def:diagonal-adaptive-estimator}\) defines \(\widehat V_n^{\mathrm{DAT}}\) directly from the unrestricted cell counts \(\overline N_x\), arm counts \(\overline N_{ax}\), and success counts \(\overline S_{ax}\), using the explicit convention \(0/0=0\).  Its formula imposes neither \(\pi_x=1/2\) nor any of the four normalized-slice mass identities.  Consequently it is a computable statistic on the full four-cell observed-margin experiment, including laws outside the embedded slice.  This proves every assertion.
\end{proof}
\par\smallskip\noindent \textit{Justification.} This proposition records the exact equal-propensity algebraic slice and prevents the lower-bound subexperiment from being confused with the full observational experiment. \textit{Gap.} The normalized slice transfers a lower bound only. Jiao--Han--Weissman Theorem 3 is a known-\(\mathbf Q\) result; simulation first carries it to the unknown--unknown comparator, while their comparator upper does not transfer to the four-cell causal model. \textit{Consumer.} A downstream reader can verify the target rescaling and the scope of the decision-theoretic reduction before importing any multinomial result.

\begin{lemma}[lem:exact-l1-transfer-parametric-floor]\label{lem:exact-l1-transfer-parametric-floor}
For every \(n\ge1\), \(d\ge2\), and \(0<\epsilon<1/2\), the embedding in \(\mathrm{def:l1-embedding}\) maps \(\Delta_d^2\) into \(\mathcal P_{d,\epsilon}\), obeys
\[
 V^\star(P)=\frac12+\frac14\lVert\mathbf P-\mathbf Q\rVert_1,
\]
and gives the exact decision-theoretic inequalities
\[
 \mathfrak R_{n,d,\epsilon}\ge\frac1{16}\mathfrak R^{L_1}_{n,d}
 \qquad\text{and}\qquad
 \mathfrak R_{n,d,\epsilon}\ge\frac{1}{1024n}.
\]
\end{lemma}
\begin{proof}
\textbf{Proof.}
For \((\mathbf P,\mathbf Q)\in\Delta_d^2\), the construction gives nonnegative four-cell masses with total mass
\[
 \sum_{x=1}^d\sum_{a,y}q_{ay,x}
 =\frac12\sum_{x=1}^d(P_x+Q_x)=1.
\]
Moreover,
\[
 p_x=\sum_{a,y}q_{ay,x}=\frac{P_x+Q_x}{2},
 \qquad
 \sum_xp_x=1,
\]
and, whenever \(p_x>0\),
\[
 \pi_x=\frac{q_{10,x}+q_{11,x}}{p_x}=\frac12,
 \quad
 \mu_{1x}=\frac{P_x}{P_x+Q_x},
 \quad
 \mu_{0x}=\frac{Q_x}{P_x+Q_x}.
\]
The potential-outcome completion in \(\mathrm{def:l1-embedding}\) satisfies consistency and conditional exchangeability by construction, while \(\pi_x=1/2\) satisfies fixed overlap because \(0<\epsilon<1/2\).  Thus the law is in \(\mathcal P_{d,\epsilon}\).  The identification formula gives
\[
 \begin{aligned}
 V^\star(P)
 &=\frac12\sum_{x=1}^d\max\{P_x,Q_x\}\\
 &=\frac14\sum_{x=1}^d\{P_x+Q_x+\lvert P_x-Q_x\rvert\}\\
 &=\frac12+\frac14\lVert\mathbf P-\mathbf Q\rVert_1.
 \end{aligned}
\]

For each \(i\), the kernel first chooses the \(\mathbf P\)-stream or the \(\mathbf Q\)-stream with probability \(1/2\), and then chooses one of the two associated \((A,Y)\) pairs with probability \(1/2\).  Hence its one-observation output probabilities are exactly
\[
 q_{11,x}=q_{00,x}=P_x/4,
 \qquad
 q_{10,x}=q_{01,x}=Q_x/4.
\]
The input streams and all kernel coins are independent across \(i\), so the output is an independent sample of size \(n\) from the embedded causal law.

Let \(\widehat V\) be any causal estimator and form the randomized comparator estimator
\[
 \widehat L=4\widehat V\circ\mathsf K_{n,d}-2.
\]
The displayed target identity implies, for every \((\mathbf P,\mathbf Q)\),
\[
 E_{\mathbf P,\mathbf Q}\!\left[
   (\widehat L-\lVert\mathbf P-\mathbf Q\rVert_1)^2
 \right]
 =16E_P[(\widehat V-V^\star(P))^2].
\]
If randomized estimators are not included in the definition of \(\mathfrak R^{L_1}_{n,d}\), replace \(\widehat L\) by its conditional expectation given the two input streams; conditional Jensen's inequality cannot increase squared risk.  Taking suprema and then the infimum over \(\widehat V\) yields
\[
 \mathfrak R^{L_1}_{n,d}\le16\mathfrak R_{n,d,\epsilon}.
\]

It remains to prove the parametric floor.  Consider two laws supported on \(X=1\), with \(\pi_1=1/2\), \(\mu_{01}=1/4\), and
\[
 \mu_{11}^{(0)}=\frac12,
 \qquad
 \mu_{11}^{(1)}=\frac12+t,
 \qquad
 t=\frac{1}{8\sqrt n}.
\]
Complete both laws by conditionally independent Bernoulli potential outcomes as in \(\mathrm{def:l1-embedding}\).  Both laws belong to \(\mathcal P_{d,\epsilon}\), and their targets differ by \(t\).  Their one-observation Hellinger affinity is
\[
 r=\frac12+\frac14\{\sqrt{1+2t}+\sqrt{1-2t}\}.
\]
The elementary inequality
\(\sqrt{1+u}+\sqrt{1-u}\ge2-u^2\) for \(\lvert u\rvert\le1/2\) gives \(r\ge1-t^2\).  Product structure and Bernoulli's inequality imply
\[
 r^n\ge(1-t^2)^n\ge1-nt^2=\frac{63}{64}.
\]
For probability masses \(f,g\), Cauchy--Schwarz gives
\[
 \lVert f-g\rVert_{\mathrm{TV}}
 \le\sqrt{1-\left(\sum_z\sqrt{f(z)g(z)}\right)^2}.
\]
Thus the total variation distance between the two \(n\)-sample laws is less than \(1/2\).  Given any estimator, test the two laws by choosing the target value closest to its output.  On a testing error the squared estimation error is at least \(t^2/4\), while the sum of the two testing-error probabilities is at least one minus total variation.  Hence the maximum of the two risks is at least
\[
 \frac{t^2}{8}\{1-\lVert P_0^n-P_1^n\rVert_{\mathrm{TV}}\}
 \ge\frac{t^2}{16}
 =\frac{1}{1024n}.
\]
Taking the infimum over estimators proves the second inequality.
\end{proof}

\begin{lemma}[lem:empirical-cellwise-ratio-upper]\label{lem:empirical-cellwise-ratio-upper}
For every \(n\ge1\), \(d\ge1\), and \(0<\epsilon<1/2\), the empirical cellwise ratio estimator in \(\mathrm{def:diagonal-adaptive-estimator}\) satisfies
\[
 \sup_{P\in\mathcal P_{d,\epsilon}}
 E_P\!\left[
  \left\{\widehat V_n^{\mathrm{DAT}}-V^\star(P)\right\}^2
 \right]
 \le
 \min\!\left\{1,
  \left(\frac12+\frac6\epsilon\right)\frac dn
 \right\}.
\]
All empirical divisions use the stated zero-denominator convention; fixed overlap is required only on cells with positive mass.
\end{lemma}
\begin{proof}
\textbf{Proof.}
Fix \(P\in\mathcal P_{d,\epsilon}\).  If \(p_x=0\), then \(\overline N_x=0\) almost surely, so cell \(x\) contributes zero to both the estimator and the target; for definiteness set \(\mu_{0x}=\mu_{1x}=0\) on such a null cell.  On every supported cell put
\[
 v_x=\max_{a\in\{0,1\}}\mu_{ax},
 \qquad
 \widehat v_x=\max_{a\in\{0,1\}}\widehat\mu_{ax}.
\]
By \(\mathrm{def:oracle-value}\) and \(\mathrm{def:diagonal-adaptive-estimator}\),
\[
 \widehat V_n^{\mathrm{DAT}}-V^\star(P)=T_1+T_2,
\]
where
\[
 T_1=\sum_{x=1}^d(\widehat p_x-p_x)v_x,
 \qquad
 T_2=\sum_{x=1}^d\widehat p_x(\widehat v_x-v_x).
\]
The first term is the centered sample average
\[
 T_1=\frac1n\sum_{i=1}^n
 \left[v_{X_i}-E_P\{v_X\}\right].
\]
Because \(0\le v_X\le1\), independent sampling in \(\mathrm{ass:iid-sampling}\) gives
\[
 E_P(T_1^2)=\frac{\operatorname{Var}_P(v_X)}n\le\frac1{4n}.
\tag{1}
\]

We next control the empirical ratios, including zero arm counts.  Fix a supported cell and condition on \(\overline N_x=k\ge1\).  Write \(\pi_{1x}=\pi_x\) and \(\pi_{0x}=1-\pi_x\).  Then
\(\overline N_{ax}\mid\overline N_x=k\) is binomial with parameters \((k,\pi_{ax})\).  Conditional on \(\overline N_{ax}=j>0\), the corresponding \(j\) outcomes are independent Bernoulli variables with mean \(\mu_{ax}\).  The convention \(0/0=0\) therefore implies
\[
 \begin{aligned}
 E_P\!\left[
   (\widehat\mu_{ax}-\mu_{ax})^2
   \mid \overline N_x=k
 \right]
 &\le \frac14 E\!\left[
    \frac{\mathbf 1\{\overline N_{ax}>0\}}{\overline N_{ax}}
    \mathrel{\Big|}\overline N_x=k
  \right]
  +P(\overline N_{ax}=0\mid\overline N_x=k)\\
 &\le \frac1{2k\pi_{ax}}+e^{-k\pi_{ax}}\\
 &\le \frac3{2k\pi_{ax}}.
\tag{2}
 \end{aligned}
\]
For the second line, use \(j^{-1}\le2(j+1)^{-1}\) for \(j\ge1\) and the binomial identity
\[
 E\!\left[
   \frac1{\overline N_{ax}+1}
   \mathrel{\Big|}\overline N_x=k
 \right]
 =\frac{1-(1-\pi_{ax})^{k+1}}{(k+1)\pi_{ax}}
 \le\frac1{k\pi_{ax}}.
\]
For the third line, use \(e^{-t}\le t^{-1}\) for \(t>0\).  Fixed overlap supplies \(\pi_{ax}\ge\epsilon\) for both arms on every supported cell; no division or overlap condition is used on null cells.

The maximum map is one-Lipschitz in the supremum norm.  Since \(\sum_x\widehat p_x=1\), weighted Cauchy--Schwarz followed by (2) gives
\[
 \begin{aligned}
 E_P(T_2^2)
 &\le E_P\!\left[
   \sum_{x=1}^d\widehat p_x(\widehat v_x-v_x)^2
 \right]\\
 &\le \sum_{x=1}^d E_P\!\left[
   \frac{\overline N_x}{n}
   \sum_{a=0}^1(\widehat\mu_{ax}-\mu_{ax})^2
 \right]\\
 &\le \frac{3d}{n\epsilon}.
\tag{3}
 \end{aligned}
\]
Indeed, on \(\{\overline N_x=k\ge1\}\), multiplication by \(k/n\) cancels the factor \(k^{-1}\), and
\[
 \frac3{2n}\left(\frac1{\pi_{0x}}+\frac1{\pi_{1x}}\right)
 \le\frac3{n\epsilon}.
\]
On \(\{\overline N_x=0\}\), the weighted summand is zero.

Equations (1)--(3) and \((u+v)^2\le2u^2+2v^2\) yield
\[
 E_P\!\left[
  \left\{\widehat V_n^{\mathrm{DAT}}-V^\star(P)\right\}^2
 \right]
 \le \frac1{2n}+\frac{6d}{n\epsilon}
 \le \left(\frac12+\frac6\epsilon\right)\frac dn,
\]
where the last inequality uses \(d\ge1\).  Finally, \(\widehat V_n^{\mathrm{DAT}}\in[0,1]\) because it is a convex combination of numbers in \([0,1]\), and \(V^\star(P)\in[0,1]\).  Thus the squared loss is at most one.  Taking the supremum over \(P\in\mathcal P_{d,\epsilon}\) proves the stated truncated bound.
\end{proof}
\par\smallskip\noindent \textit{Justification.} This one-pass estimator is the strongest fully certified achievability result under the unchanged unrestricted observational model and handles null cells and zero arm counts explicitly. \textit{Gap.} Its \(d/n\) risk does not use a heavy/light polynomial gain. The accepted linear-ATE coefficient certificate controls linear signed ratios, but no simultaneous approximation and absolute-coefficient certificate is proved for the nonlinear \(\Gamma\) term. \textit{Consumer.} A practitioner can compute the estimator directly from cell counts and invoke a uniform finite-sample MSE guarantee without tuning or a pilot split.

\begin{lemma}[lem:alphabet-monotonicity]\label{lem:alphabet-monotonicity}
For every \(n\geq1\), \(d\geq m\geq2\), and \(0<\epsilon<1/2\), null-cell padding embeds \(\mathcal P_{m,\epsilon}\) into \(\mathcal P_{d,\epsilon}\): each embedded law has \(p_x=0\) for \(x>m\), has the same observed sample law and the same value \(V^\star(P)\) as the original law, and may use arbitrary observationally irrelevant conditional coordinates on the null cells. Consequently,
\[
 \mathfrak R_{n,d,\epsilon}\geq \mathfrak R_{n,m,\epsilon}.
\]
\end{lemma}
\begin{proof}
\textbf{Proof.}
Fix \(n\geq1\), \(d\geq m\geq2\), and \(0<\epsilon<1/2\).  Let \(P\in\mathcal P_{m,\epsilon}\).  Regard the same random variables as taking values in the larger alphabet \([d]\), under the canonical inclusion \([m]\subset[d]\), and denote the resulting law by \(P^{\uparrow d}\).  Then
\[
 p_x^{\uparrow d}=p_x\quad(1\leq x\leq m),
 \qquad
 p_x^{\uparrow d}=0\quad(m<x\leq d).
\]
The consistency and conditional-exchangeability statements are unchanged because the joint law of \((X,A,Y,Y(0),Y(1))\) is unchanged.  Fixed overlap continues to hold on every supported cell.  It imposes no condition on the newly added cells because their masses are zero.  Thus \(P^{\uparrow d}\in\mathcal P_{d,\epsilon}\).  Conditional coordinates on the null cells can, for example, be represented by propensity \(1/2\) and deterministic zero potential outcomes; they do not change the law because those cells have probability zero.  This verifies the support, positivity, overlap, and boundary conditions of the embedding.

By \(\mathrm{def:oracle-value}\), all newly added summands vanish, so
\[
 V^\star(P^{\uparrow d})
 =\sum_{x=1}^m p_x\max\{\mu_{0x},\mu_{1x}\}
 =V^\star(P).
\]
The observed \(n\)-sample law under \(P^{\uparrow d}\), after the canonical inclusion of its sample space, is exactly the observed \(n\)-sample law under \(P\).

It remains to check the direction of the decision-space comparison.  Let \(T_d\) be any measurable estimator on \(([d]\times\{0,1\}^2)^n\).  Its restriction
\[
 T_m(o_1,\ldots,o_n):=T_d(o_1,\ldots,o_n),
 \qquad o_i\in[m]\times\{0,1\}^2,
\]
is a measurable estimator on the \(m\)-alphabet experiment.  Conversely, every measurable \(T_m\) has a measurable extension to the \(d\)-alphabet sample space, for instance by setting the extension equal to zero whenever at least one observed covariate lies outside \([m]\).  Hence the restrictions of the \(d\)-alphabet decision rules are exactly all \(m\)-alphabet decision rules.  Since the padded laws form a subclass of \(\mathcal P_{d,\epsilon}\), \(\mathrm{def:minimax-risk}\) gives
\[
\begin{aligned}
 \mathfrak R_{n,d,\epsilon}
 &=\inf_{T_d}\sup_{Q\in\mathcal P_{d,\epsilon}}
   E_Q\!\left[\{T_d-V^\star(Q)\}^2\right]\\
 &\geq\inf_{T_d}\sup_{P\in\mathcal P_{m,\epsilon}}
   E_{P^{\uparrow d}}\!\left[\{T_d-V^\star(P^{\uparrow d})\}^2\right]\\
 &=\inf_{T_m}\sup_{P\in\mathcal P_{m,\epsilon}}
   E_P\!\left[\{T_m-V^\star(P)\}^2\right]
 =\mathfrak R_{n,m,\epsilon}.
\end{aligned}
\]
This proves the claimed alphabet monotonicity.
\end{proof}
\par\smallskip\noindent \textit{Justification.} This lemma is the maximality step: padding with zero-mass cells preserves membership, the observed experiment, and the causal value, while the estimator-restriction argument fixes the minimax inequality direction. \textit{Gap.} The earlier range-limited lower bounds alone did not control alphabets above their construction range. \textit{Consumer.} The global bracket truncates the hard alphabet at order \(n\log(en)\) and obtains a constant lower floor for every larger alphabet.

\begin{lemma}[lem:jiao-han-weissman-known-q-lower]\label{lem:jiao-han-weissman-known-q-lower}
Fix \(C>0\).  In the Poissonized known-\(\mathbf Q\) experiment of Jiao, Han, and Weissman, let
\[
 T_{n,d}(\mathbf Q)
 :=\sum_{x=1}^d\left\{Q_x\wedge\sqrt{\frac{Q_x}{n\log n}}\right\}.
\]
Their Theorem 3 states that, when \(d\ge2\), \(\log n\ge C\log d\), and
\[
 T_{n,d}(\mathbf Q)
 \ge C'\left\{\sqrt{\frac{\log n}{n}}+\frac{\sqrt d\log n}{n}\right\},
\]
where \(C'>0\) depends only on \(C\), the minimax mean-squared error for estimating \(\lVert\mathbf P-\mathbf Q\rVert_1\) is bounded below by \(c_C T_{n,d}(\mathbf Q)^2\).  The Poissonized-to-fixed-multinomial correction is supplied by Jiao, Han, and Weissman, Lemma 33 (proof, equations (440)--(459), arXiv:1705.00807v7).
\end{lemma}
\par\smallskip\noindent \textit{Justification.} This cited leaf isolates the published known-\(\mathbf Q\) lower bound and the exact fixed-sample correction used by the high-alphabet comparator argument. \textit{Gap.} Theorem 3 is a known-\(\mathbf Q\) result, and the Poissonized-to-fixed-multinomial attribution is specifically Lemma 33 of arXiv:1705.00807v7. Simulation transfers the lower bound to the unknown--unknown comparator, but neither cited result supplies an upper bound for the observational four-cell causal experiment. \textit{Consumer.} The two-band normalized-\(L_1\) lower lemma applies the cited bound at uniform \(\mathbf Q\) and then the exact causal embedding transfers it to \(\mathfrak R_{n,d,\epsilon}\).

\begin{lemma}[lem:two-band-l1-lower]\label{lem:two-band-l1-lower}
There are universal constants \(c,C,\rho>0\) such that the two-sample multinomial risk in \(\mathrm{def:l1-embedding}\) satisfies
\[
 \mathfrak R^{L_1}_{n,d}\ge c\frac{d}{n\log(en)}
\]
for every \(2\le d\le\rho n\log(en)\).  For \(d\le C\log^2(en)\) this follows from a paired-sign fuzzy mixture, and for \(d\ge C\log^2(en)\) it follows from \(\mathrm{lem:jiao-han-weissman-known-q-lower}\) with uniform \(\mathbf Q\).
\end{lemma}
\begin{proof}
\textbf{Proof.}
Write \(L=\log(en)=1+\log n\).  We first prove the lower-dimensional band directly and then invoke the cited known-\(\mathbf Q\) theorem in the complementary band.

\emph{Paired-sign band.}
Let \(k=2\lfloor d/2\rfloor\), so that \(k\) is even and \(d/2\le k\le d\).  Define \(\mathbf Q^{(0)}\in\Delta_d\) by \(Q^{(0)}_x=1/k\) for \(x\le k\) and \(Q^{(0)}_x=0\) otherwise.  For \(\sigma=(\sigma_1,\ldots,\sigma_{k/2})\in\{-1,1\}^{k/2}\) and \(0<\alpha\le1/2\), define \(\mathbf P_\sigma\in\Delta_d\) by
\[
 P_{\sigma,2j-1}=\frac{1+\alpha\sigma_j}{k},\qquad
 P_{\sigma,2j}=\frac{1-\alpha\sigma_j}{k},
 \qquad 1\le j\le k/2,
\]
and set its remaining coordinates to zero.  Nonnegativity follows from \(\alpha\le1/2\), and summing within pairs shows that the coordinates sum to one.  Furthermore,
\[
 \lVert\mathbf P_\sigma-\mathbf Q^{(0)}\rVert_1
 =\sum_{j=1}^{k/2}\left(\frac{\alpha}{k}+\frac{\alpha}{k}\right)
 =\alpha .
\]

Let \(M\) be the uniform mixture of the laws \(\mathbf P_\sigma^{\otimes n}\), and let \(Q_0=(\mathbf Q^{(0)})^{\otimes n}\).  If \(Z\sim\mathbf Q^{(0)}\), direct multiplication of the one-observation likelihood ratios gives, for every \(\sigma,\tau\),
\[
 E\left[
 \frac{P_\sigma(Z)}{Q^{(0)}(Z)}
 \frac{P_\tau(Z)}{Q^{(0)}(Z)}
 \right]
 =1+\frac{2\alpha^2}{k}\sum_{j=1}^{k/2}\sigma_j\tau_j.
\]
Independence of the \(n\) observations and expansion of the squared mixture likelihood ratio therefore yield
\[
 1+\chi^2(M,Q_0)
 =E_{\sigma,\tau}\left(
 1+\frac{2\alpha^2}{k}\sum_{j=1}^{k/2}\sigma_j\tau_j
 \right)^n.
\]
Using \(1+t\le e^t\), independence of the Rademacher products \(\sigma_j\tau_j\), and \(\cosh t\le e^{t^2/2}\), we obtain
\[
 1+\chi^2(M,Q_0)
 \le
 \left\{\cosh\left(\frac{2n\alpha^2}{k}\right)\right\}^{k/2}
 \le \exp\left(\frac{n^2\alpha^4}{k}\right).
\]
Choose a universal \(b>0\) sufficiently small that
\[
 \frac12\sqrt{e^{b^2}-1}\le\frac12,
\]
and set \(\alpha^2=b\sqrt{k}/n\).  Whenever \(d\le K L^2\), for any fixed universal \(K<\infty\), this choice satisfies \(\alpha\le1/2\) once \(n\) exceeds a universal threshold depending only on \(K\).  The preceding display then gives \(\chi^2(M,Q_0)\le e^{b^2}-1\).  Cauchy--Schwarz applied to the mixture likelihood ratio gives
\[
 \lVert M-Q_0\rVert_{\mathrm{TV}}
 \le\frac12\sqrt{\chi^2(M,Q_0)}\le\frac12.
\]

We spell out the risk reduction.  For an arbitrary estimator \(\widehat L\), let \(\phi=\mathbf 1\{\widehat L\ge\alpha/2\}\).  Under the null pair \((\mathbf Q^{(0)},\mathbf Q^{(0)})\), the target is zero, whereas under every alternative pair \((\mathbf P_\sigma,\mathbf Q^{(0)})\), it is \(\alpha\).  The second sample has the same law under the null and all alternatives, so tensoring both \(M\) and \(Q_0\) with that common law leaves their total-variation distance unchanged.  Consequently,
\[
 \begin{aligned}
 &\max\left\{
 E_{\mathbf Q^{(0)},\mathbf Q^{(0)}}\widehat L^2,
 \sup_\sigma E_{\mathbf P_\sigma,\mathbf Q^{(0)}}(\widehat L-\alpha)^2
 \right\}\\
 &\quad\ge
 \frac{\alpha^2}{8}
 \left\{Q_0(\phi=1)+M(\phi=0)\right\}
 \ge \frac{\alpha^2}{8}
 \left\{1-\lVert M-Q_0\rVert_{\mathrm{TV}}\right\}
 \ge\frac{\alpha^2}{16}.
 \end{aligned}
\]
Taking the infimum over \(\widehat L\) proves
\[
 \mathfrak R^{L_1}_{n,d}\ge c\frac{\sqrt{k}}{n}.
\]
If \(d\le K L^2\), then \(k\ge d/2\) and hence
\[
 \frac{\sqrt{k}}n
 \ge \frac{d}{\sqrt{2K}\,nL}.
\]
Thus the asserted lower bound holds throughout this band.

\emph{Uniform known-\(\mathbf Q\) band.}
Apply \(\mathrm{lem:jiao-han-weissman-known-q-lower}\) with its logarithmic constant fixed at \(1/2\), and denote the corresponding constants by \(C'>0\) and \(c'>0\).  Take \(\mathbf Q^{(u)}=(1/d,\ldots,1/d)\).  We will choose \(\rho>0\) small enough that, for all sufficiently large \(n\),
\[
 d\le\rho nL\quad\Longrightarrow\quad d\le n\log n.
\]
For such \(n,d\), the minimum in the cited theorem is attained by its square-root term, and therefore
\[
 T_{n,d}(\mathbf Q^{(u)})
 =d\sqrt{\frac{1/d}{n\log n}}
 =\sqrt{\frac d{n\log n}}.
\]
Choose the universal constant \(K\) large enough that \(d\ge K L^2\) implies
\[
 T_{n,d}(\mathbf Q^{(u)})
 \ge C'\sqrt{\frac{\log n}{n}}.
\]
Moreover, identically in \(d\),
\[
 \frac{T_{n,d}(\mathbf Q^{(u)})}{\sqrt d\log n/n}
 =\frac{\sqrt n}{(\log n)^{3/2}},
\]
which is at least \(C'\) for all sufficiently large \(n\).  Finally, if \(d\le\rho nL\), then \(\log d\le2\log n\) for all sufficiently large \(n\), so the cited requirement \(\log n\ge\frac12\log d\) also holds.  All hypotheses of the cited lower bound are therefore satisfied when \(d\ge K L^2\), and its fixed-multinomial comparison gives
\[
 \inf_{\widehat L}
 \sup_{\mathbf P\in\Delta_d}
 E_{\mathbf P,\mathbf Q^{(u)}}
 \left(\widehat L-\lVert\mathbf P-\mathbf Q^{(u)}\rVert_1\right)^2
 \ge c'\frac d{n\log n}
 \ge c'\frac d{nL}.
\]
The unrestricted two-sample minimax risk cannot be smaller: restricting its supremum to \(\mathbf Q=\mathbf Q^{(u)}\) and then revealing this fixed \(\mathbf Q^{(u)}\) only makes the estimation problem easier.  Hence the same lower bound holds for \(\mathfrak R^{L_1}_{n,d}\).

It remains only to choose constants uniformly.  First choose \(K\) as above, and then choose a universal integer \(n_0\) large enough for every displayed large-\(n\) condition, including \(b\sqrt{k}/n\le1/4\) when \(k\le K L^2\).  Choose \(\rho>0\) so small that the uniform-band implications hold for \(n\ge n_0\) and, in addition,
\[
 \rho n\log(en)<2\qquad(1\le n<n_0).
\]
For \(n<n_0\) the claimed range \(2\le d\le\rho n\log(en)\) is then empty.  For \(n\ge n_0\), either \(d\le K L^2\) and the paired-sign argument applies, or \(d\ge K L^2\) and the cited uniform argument applies.  Decreasing \(c>0\) to the smaller of the two universal band constants and taking \(C=K\) proves the result. \(\square\)
\end{proof}
\par\smallskip\noindent \textit{Justification.} This helper makes the normalized-distance converse valid over the full range by joining a proved fuzzy-mixture band to the precisely cited known-\(\mathbf Q\) band. \textit{Gap.} The cited theorem's dominance condition begins only above a logarithmic-square threshold, so the lower band must be proved directly. \textit{Consumer.} The exact causal reduction consumes this helper to obtain the \(d/[n\log(en)]\) lower term.

\begin{theorem}[thm:elementary-minimax-bracket]\label{thm:elementary-minimax-bracket}
For every fixed \(0<\epsilon<1/2\), there are constants \(0<c_\epsilon\le C_\epsilon<\infty\) and \(\rho_\epsilon>0\) such that, uniformly for \(2\le d\le\rho_\epsilon n\log(en)\),
\[
 c_\epsilon\left\{\frac{\sqrt d}{n}+\frac d{n\log(en)}\right\}
 \le\mathfrak R_{n,d,\epsilon}
 \le\sup_{P\in\mathcal P_{d,\epsilon}}
 E_P\!\left[
  \left\{\widehat V_n^{\mathrm{DAT}}-V^\star(P)\right\}^2
 \right]
 \le\min\!\left\{1,C_\epsilon\frac dn\right\}.
\]
The two displayed lower-rate terms balance at \(d\asymp\{\log(en)\}^2\): the paired-sign term \(\sqrt d/n\) dominates below this elbow and the normalized-distance term \(d/\{n\log(en)\}\) dominates above it.  Along sequences in the displayed range, \(d=o(n\log n)\) is necessary for consistency and \(d=o(n)\) is sufficient.  Moreover, \(\mathfrak R_{n,d,\epsilon}=\Theta(n^{-1})\) if and only if \(d=O(1)\).  The bracket does not close the intermediate consistency region or matching achievability in either lower-rate regime.
\end{theorem}
\begin{proof}
\textbf{Proof.}
Apply \(\mathrm{thm:global-elementary-minimax-bracket}\).  Its upper bound is already the asserted one.  For the lower bound, choose the range constant \(\rho_\epsilon>0\) no larger than \(1/4\).  If \(d\le\rho_\epsilon n\log(en)\), then
\[
 \frac d{n\log(en)}\le\rho_\epsilon,
 \qquad
 \frac{\sqrt d}{n}
 \le\sqrt{\rho_\epsilon\frac{\log(en)}n}
 \le\sqrt{\rho_\epsilon},
\]
where \(\log(en)/n\le1\) for every integer \(n\ge1\).  Hence the sum inside the global truncation is at most \(\sqrt{\rho_\epsilon}+\rho_\epsilon<1\), so the truncation disappears.  This gives the displayed in-range lower bound after retaining the same \(c_\epsilon\).

The comparison
\[
 \frac{\sqrt d/n}{d/\{n\log(en)\}}
 =\frac{\log(en)}{\sqrt d}
\]
proves the elbow and dominance assertions.  The necessary and sufficient consistency statements, the parametric iff, and the explicit description of the remaining gap are restrictions of the corresponding conclusions of the global theorem to the displayed range.  No additional support, positivity, or boundary condition is introduced by this restriction.
\end{proof}
\par\smallskip\noindent \textit{Justification.} The retained range statement is an uncapped corollary of the global padded theorem and preserves existing references. \textit{Gap.} Its former range qualifier is no longer a limitation of the available lower bound; the global theorem is the canonical strengthened result. \textit{Consumer.} Existing consumers of the range-limited node receive the same finite-sample inequalities with a shorter proof spine.

\begin{lemma}[lem:paired-sign-small-alphabet-floor]\label{lem:paired-sign-small-alphabet-floor}
For every \(n,d\in\mathbb N\), \(0<\epsilon<1/2\), and \(2\le d\le n^2\), put \(k=2\lfloor d/2\rfloor\).  Let \(\mathbf Q\in\Delta_d\) be uniform on its first \(k\) coordinates and zero on the rest.  For \(\sigma\in\{-1,1\}^{k/2}\), define \(\mathbf P_\sigma\in\Delta_d\) by
\[
 P_{\sigma,2j-1}=\frac{1+\alpha\sigma_j}{k},\qquad
 P_{\sigma,2j}=\frac{1-\alpha\sigma_j}{k},
 \qquad
 \alpha^2=\frac{\sqrt{k}}{16n},
\]
and set its remaining coordinates to zero.  The uniform mixture of the \(n\)-sample laws \(\mathbf P_\sigma^n\) has chi-square divergence at most \(e^{1/256}-1\) from \(\mathbf Q^n\), while \(\lVert\mathbf P_\sigma-\mathbf Q\rVert_1=\alpha\) for every \(\sigma\).  Consequently,
\[
 \mathfrak R^{L_1}_{n,d}\ge \frac{\sqrt{k}}{256n}
 \qquad\text{and}\qquad
 \mathfrak R_{n,d,\epsilon}
 \ge \frac{\sqrt d}{4096\sqrt2\,n}.
\]
\end{lemma}
\begin{proof}
\textbf{Proof.}
Because \(2\le d\le n^2\) and \(k=2\lfloor d/2\rfloor\le d\),
\[
 0<\alpha^2=\frac{\sqrt{k}}{16n}\le\frac1{16}.
\]
Thus all displayed coordinates of \(\mathbf P_\sigma\) are nonnegative, every pair sums to \(2/k\), and \(\mathbf P_\sigma\in\Delta_d\).  Coordinatewise comparison with the stated uniform \(\mathbf Q\) gives
\[
 \lVert\mathbf P_\sigma-\mathbf Q\rVert_1
 =\sum_{j=1}^{k/2}\left(\frac\alpha k+\frac\alpha k\right)
 =\alpha.
\]

Let \(M=2^{-k/2}\sum_\sigma\mathbf P_\sigma^n\).  If \(L_\sigma=d\mathbf P_\sigma^n/d\mathbf Q^n\), then independence and the finite likelihood-ratio identity yield
\[
 \begin{aligned}
 1+\chi^2(M,\mathbf Q^n)
 &=E_{\sigma,\tau}E_{\mathbf Q^n}(L_\sigma L_\tau)\\
 &=E_{\sigma,\tau}\left(
    \sum_{i=1}^k\frac{P_{\sigma,i}P_{\tau,i}}{Q_i}
   \right)^n\\
 &=E_{\sigma,\tau}\left(
    1+\frac{2\alpha^2}{k}
      \sum_{j=1}^{k/2}\sigma_j\tau_j
   \right)^n.
 \end{aligned}
\]
The last base is at least \(1-\alpha^2>0\).  Setting \(\xi_j=\sigma_j\tau_j\), using \(1+u\le e^u\), and then using \(\cosh t\le e^{t^2/2}\), we obtain
\[
 \begin{aligned}
 1+\chi^2(M,\mathbf Q^n)
 &\le E\exp\left{
   \frac{2n\alpha^2}{k}\sum_{j=1}^{k/2}\xi_j
  \right}\\
 &=\left\{\cosh\left(\frac{2n\alpha^2}{k}\right)\right\}^{k/2}
 \le\exp\left(\frac{n^2\alpha^4}{k}\right)
 =e^{1/256}.
\end{aligned}
\]
This proves the asserted chi-square bound.  The independent second sample from \(\mathbf Q\) is common to the null and the mixture, so it changes neither chi-square divergence nor total variation.

Now take any estimator \(\widehat L\) in the two-sample experiment and test the null \((\mathbf Q,\mathbf Q)\) against the fuzzy alternative by rejecting when \(\widehat L\ge\alpha/2\).  On either type of testing error the squared estimation loss is at least \(\alpha^2/4\).  Since
\[
 \lVert M-\mathbf Q^n\rVert_{\mathrm{TV}}
 \le\frac12\sqrt{\chi^2(M,\mathbf Q^n)}
 \le\frac12\sqrt{e^{1/256}-1}<\frac1{16},
\]
the sum of the null error and mixture-average alternative error is at least \(15/16\).  The larger of the null risk and the mixture-average alternative risk is consequently at least
\[
 \frac{\alpha^2}{8}\frac{15}{16}
 \ge\frac{\alpha^2}{16}
 =\frac{\sqrt{k}}{256n}.
\]
Taking the infimum over estimators proves the comparator lower bound.  The exact transfer in \(\mathrm{lem:l1-lower-transfer}\) then gives
\[
 \mathfrak R_{n,d,\epsilon}
 \ge\frac{\sqrt{k}}{4096n}.
\]
Finally \(k=2\lfloor d/2\rfloor\ge d/2\) for \(d\ge2\), so
\[
 \mathfrak R_{n,d,\epsilon}
 \ge\frac{\sqrt d}{4096\sqrt2\,n}.
\]
All alternative laws used by the transfer have propensity \(1/2\), which satisfies the fixed-overlap boundary \(0<\epsilon<1/2\).
\end{proof}

\begin{theorem}[thm:no-four-cell-logarithmic-upper]\label{thm:no-four-cell-logarithmic-upper}
For every fixed \(0<\epsilon<1/2\), there do not exist finite \(C_\epsilon\), positive \(\rho_\epsilon\), and a collection of measurable estimators \(\widehat V_{n,d}\) such that
\[
 \sup_{P\in\mathcal P_{d,\epsilon}}
 E_P\!\left[
  \left\{\widehat V_{n,d}(O_1,\ldots,O_n)-V^\star(P)\right\}^2
 \right]
 \le C_\epsilon\left\{\frac1n+\frac d{n\log(en)}\right\}
\]
for every \(2\le d\le\rho_\epsilon n\log(en)\).  Indeed, for \(d_n=\max\{2,\lfloor\log(en)\rfloor\}\),
\[
 \mathfrak R_{n,d_n,\epsilon}
 \ge \frac{\sqrt{\log(en)}}{8192n}
\]
for all sufficiently large \(n\).  Hence no finite computable four-cell overlap-cone construction can have the proposed uniform risk certificate.  Moreover, along every sequence with \(2\le d\le n^2\), \(\mathfrak R_{n,d,\epsilon}=O(n^{-1})\) holds if and only if \(d=O(1)\); for bounded \(d\), the risk is of order \(n^{-1}\).
\end{theorem}
\begin{proof}
\textbf{Proof.}
Let \(L_n=\log(en)\) and \(d_n=\max\{2,\lfloor L_n\rfloor\}\).  For all sufficiently large \(n\),
\[
 2\le d_n\le n^2,
 \qquad
 d_n\ge\frac{L_n}{2}.
\]
Lemma \(\mathrm{lem:paired-sign-small-alphabet-floor}\) therefore gives
\[
 \mathfrak R_{n,d_n,\epsilon}
 \ge\frac{\sqrt{d_n}}{4096\sqrt2\,n}
 \ge\frac{\sqrt{L_n}}{8192n}.
\tag{11}
\]

Suppose, toward a contradiction, that finite \(C_\epsilon\), positive \(\rho_\epsilon\), and estimators satisfying the displayed uniform certificate existed.  Since
\[
 \frac{d_n}{nL_n}\longrightarrow0,
\]
the point \(d=d_n\) belongs to \(2\le d\le\rho_\epsilon nL_n\) for every sufficiently large \(n\).  By \(\mathrm{def:minimax-risk}\), the assumed estimator bound would imply
\[
 \mathfrak R_{n,d_n,\epsilon}
 \le C_\epsilon\left\{\frac1n+\frac{d_n}{nL_n}\right\}
 \le\frac{2C_\epsilon}{n}.
\tag{12}
\]
Equations (11) and (12) contradict one another because \(\sqrt{L_n}\to\infty\).  Any finite computable four-cell construction carrying the proposed certificate would define a measurable estimator collection of the prohibited kind, so the same contradiction excludes that certificate.

For the final assertion, take any sequence satisfying \(2\le d\le n^2\).  If \(d=O(1)\), \(\mathrm{thm:global-elementary-minimax-bracket}\) gives an \(O(n^{-1})\) upper bound and its parametric floor gives a matching \(\Omega(n^{-1})\) lower bound.  Conversely, the paired-sign lemma gives
\[
 \mathfrak R_{n,d,\epsilon}
 \ge\frac{\sqrt d}{4096\sqrt2\,n}.
\]
Thus an \(O(n^{-1})\) risk bound forces \(\sqrt d=O(1)\), equivalently \(d=O(1)\).  This proves both directions.
\end{proof}
\par\smallskip\noindent \textit{Justification.} The new obstruction is a concrete causal sequence witness at \(d_n\asymp\log(en)\): the paired-sign mixture forces risk \(\sqrt{\log(en)}/n\), ruling out the proposed direct four-cell lift of the split-sample diagonal-adaptive workflow in JiaoHanWeissman2018L1, Construction 2 and Theorem 6, with the advertised \(n^{-1}+d/[n\log(en)]\) certificate. \textit{Gap.} JiaoHanWeissman2018L1, Construction 2 and Theorem 6, prove an unknown--unknown normalized-multinomial \(L_1\) estimator with \(d/[n\log n]\)-scale risk in that experiment; the present theorem does not dispute that result. It disproves its direct transplantation to the larger observational four-cell ratio model at the stronger proposed certificate. Luedtke--van der Laan explains local tie nonregularity, and Whitehouse et al. recover root-\(n\) behavior under gap-density restrictions. None gives this unrestricted growing-alphabet witness. \textit{Consumer.} It prevents future estimator designs from treating a direct four-cell lift of the Jiao--Han--Weissman workflow, or the abandoned logarithmic certificate, as a valid benchmark; any successful upper theorem must pay or otherwise overcome the paired-sign cost.

\begin{theorem}[thm:no-universal-leading-constant]\label{thm:no-universal-leading-constant}
For every fixed \(0<\epsilon<1/2\), there is no finite leading constant under the normalization \(n\log n/d\) uniformly over all sequences satisfying \(d/\log n\to\infty\) and \(d/(n\log n)\to0\).  With \(L_n=\log(en)\) and \(d_n=\lfloor L_n^{3/2}\rfloor\), these two sequence conditions hold, but
\[
 \liminf_{n\to\infty}
 \frac{n\log n}{d_n}\,
 \mathfrak R_{n,d_n,\epsilon}=+\infty.
\]
Consequently, no finite universal limit over the regime in the question can be characterized or attained through \(\mathscr H_\epsilon^{\mathrm{dual}}\).  The proved lower bounds instead motivate separate unresolved exact-constant questions with normalization \(n/\sqrt d\) below the elbow, normalization \(n\log n/d\) when \(d/\log^2 n\to\infty\), and an elbow analysis when \(d\asymp\kappa\log^2 n\); no convergence or attainability assertion for those narrower regimes is made here.
\end{theorem}
\begin{proof}
\textbf{Proof.}
Put \(L_n=\log(en)\) and \(d_n=\lfloor L_n^{3/2}\rfloor\).  Then, as \(n\to\infty\),
\[
 \frac{d_n}{\log n}\sim L_n^{1/2}\longrightarrow\infty,
 \qquad
 \frac{d_n}{n\log n}\sim\frac{L_n^{1/2}}n\longrightarrow0.
\]
Also \(2\le d_n\le n^2\) for all sufficiently large \(n\).  Lemma \(\mathrm{lem:paired-sign-small-alphabet-floor}\) and the elementary relation \(d_n\sim L_n^{3/2}\) give
\[
 \begin{aligned}
 \frac{n\log n}{d_n}\,
 \mathfrak R_{n,d_n,\epsilon}
 &\ge
 \frac{\log n}{4096\sqrt2\,\sqrt{d_n}}\\
 &\sim\frac{1}{4096\sqrt2}L_n^{1/4}
 \longrightarrow+\infty.
 \end{aligned}
\]
This proves the displayed liminf and rules out a finite universal leading constant over the stated collection of sequences.

The program \(\mathscr H_\epsilon^{\mathrm{dual}}\) is only a handle for a proposed constant calculation; it cannot characterize or attain a finite constant under a normalization for which the minimax risk just proved diverges.  The paired-sign lower bound scales as \(\sqrt d/n\) below the logarithmic-square elbow, whereas the normalized-distance lower bound scales as \(d/\{n\log n\}\) above it.  Therefore the three narrower normalizations listed in the statement are the only conclusions motivated here.  No existence, convergence, or attainability claim for their constants follows from the present bounds, exactly as stated.
\end{proof}
\par\smallskip\noindent \textit{Justification.} A concrete sequence below the elbow makes the proposed \(n\log n/d\) normalization diverge. \textit{Gap.} Exact constants may only be posed separately below, above, and at the logarithmic-square elbow. \textit{Consumer.} It limits any future approximation-duality calculation to a regime with the correct normalization.

\begin{theorem}[thm:global-elementary-minimax-bracket]\label{thm:global-elementary-minimax-bracket}
For every fixed \(0<\epsilon<1/2\), there are constants \(0<c_\epsilon\leq C_\epsilon<\infty\) such that, for every \(n\geq1\) and \(d\geq2\),
\[
 c_\epsilon\min\!\left\{1,
   \frac{\sqrt d}{n}+\frac d{n\log(en)}
 \right\}
 \leq\mathfrak R_{n,d,\epsilon}
 \leq\sup_{P\in\mathcal P_{d,\epsilon}}
 E_P\!\left[
  \left\{\widehat V_n^{\mathrm{DAT}}-V^\star(P)\right\}^2
 \right]
 \leq\min\!\left\{1,C_\epsilon\frac dn\right\}.
\]
In the nonsaturated range the lower terms balance at \(d\asymp\{\log(en)\}^2\), with \(\sqrt d/n\) dominant below the elbow and \(d/\{n\log(en)\}\) dominant above it.  Along arbitrary sequences with \(d\geq2\), consistency requires \(d=o(n\log n)\), while \(d=o(n)\) is sufficient.  Moreover, \(\mathfrak R_{n,d,\epsilon}=\Theta(n^{-1})\) if and only if \(d=O(1)\).  The displayed bracket does not settle consistency in the intermediate region or matching achievability in either nonsaturated lower-rate regime.
\end{theorem}
\begin{proof}
\textbf{Proof.}
Fix \(0<\epsilon<1/2\) and write \(L_n=\log(en)\).  We first globalize the paired-sign lower bound.  By \(\mathrm{lem:paired-sign-small-alphabet-floor}\), for \(2\le d\le n^2\),
\[
 \mathfrak R_{n,d,\epsilon}\ge c_0\frac{\sqrt d}{n},
 \qquad c_0=\frac1{4096\sqrt2}.
\tag{4}
\]
If \(n\ge2\) and \(d>n^2\), apply \(\mathrm{lem:alphabet-monotonicity}\) with \(m=n^2\), and then apply (4) at alphabet size \(m\).  This gives
\[
 \mathfrak R_{n,d,\epsilon}
 \ge\mathfrak R_{n,n^2,\epsilon}\ge c_0.
\]
For \(n=1\), the parametric floor in \(\mathrm{lem:l1-lower-transfer}\) supplies a positive constant for every \(d\ge2\).  After decreasing the constant, these cases combine as
\[
 \mathfrak R_{n,d,\epsilon}
 \ge c_{1,\epsilon}\min\!\left\{1,\frac{\sqrt d}{n}\right\}.
\tag{5}
\]

Next let \(c_{2,\epsilon}>0\) and \(\rho_0>0\) be constants supplied by \(\mathrm{lem:l1-lower-transfer}\).  In the range \(2\le d\le\rho_0nL_n\), that lemma gives
\[
 \mathfrak R_{n,d,\epsilon}
 \ge c_{2,\epsilon}\frac d{nL_n}.
\tag{6}
\]
Choose \(N_0<\infty\) so that \(\rho_0nL_n\ge8\) for \(n\ge N_0\).  If \(n\ge N_0\) and \(d>\rho_0nL_n\), put
\[
 m=\left\lfloor\frac{\rho_0nL_n}{2}\right\rfloor.
\]
Then \(2\le m\le d\), \(m\le\rho_0nL_n\), and \(m\ge\rho_0nL_n/4\).  Alphabet monotonicity and (6) imply
\[
 \mathfrak R_{n,d,\epsilon}
 \ge\mathfrak R_{n,m,\epsilon}
 \ge c_{2,\epsilon}\frac m{nL_n}
 \ge\frac{c_{2,\epsilon}\rho_0}{4}.
\]
For the finitely many \(n<N_0\), the parametric floor \(1/(1024n)\) is at least \(1/(1024N_0)\).  Decreasing the constant once more gives, for every \(n\ge1\) and \(d\ge2\),
\[
 \mathfrak R_{n,d,\epsilon}
 \ge c_{3,\epsilon}\min\!\left\{1,\frac d{nL_n}\right\}.
\tag{7}
\]

For nonnegative \(u,v\),
\[
 \max\{\min(1,u),\min(1,v)\}
 \ge\frac12\{\min(1,u)+\min(1,v)\}
 \ge\frac12\min(1,u+v).
\tag{8}
\]
Apply (8) to (5) and (7), with \(u=\sqrt d/n\) and \(v=d/(nL_n)\), and decrease the lower constant if necessary.  This proves
\[
 \mathfrak R_{n,d,\epsilon}
 \ge c_\epsilon\min\!\left\{1,
   \frac{\sqrt d}{n}+\frac d{nL_n}
 \right\}.
\tag{9}
\]

For the upper bound, the infimum in \(\mathrm{def:minimax-risk}\) may be evaluated at the single data-driven estimator \(\widehat V_n^{\mathrm{DAT}}\).  Lemma \(\mathrm{lem:empirical-cellwise-ratio-upper}\) gives
\[
 \begin{aligned}
 \mathfrak R_{n,d,\epsilon}
 &\le\sup_{P\in\mathcal P_{d,\epsilon}}
 E_P\!\left[
   \left\{\widehat V_n^{\mathrm{DAT}}-V^\star(P)\right\}^2
 \right]\\
 &\le\min\!\left\{1,
  \left(\frac12+\frac6\epsilon\right)\frac dn
 \right\}.
\tag{10}
 \end{aligned}
\]
Thus one may take \(C_\epsilon=1/2+6/\epsilon\), enlarged if necessary so that \(c_\epsilon\le C_\epsilon\).

The ratio of the two nonsaturated lower terms is
\[
 \frac{\sqrt d/n}{d/(nL_n)}=\frac{L_n}{\sqrt d}.
\]
Consequently they cross at \(d\asymp L_n^2\): the first proved lower mechanism is larger below this crossing and the second is larger above it.  Because (10) is not matched to (9), this algebraic crossing is not asserted to be a matched minimax phase boundary.

Now consider sequences indexed by \(n\to\infty\).  If \(\mathfrak R_{n,d,\epsilon}\to0\), (7) forces
\[
 \frac d{nL_n}\longrightarrow0,
\]
which is equivalent to \(d=o(n\log n)\).  Conversely, \(d=o(n)\) makes (10) tend to zero.

Finally, if \(d=O(1)\), (10) is \(O(n^{-1})\), while the parametric floor in \(\mathrm{lem:l1-lower-transfer}\) is \(\Omega(n^{-1})\).  Hence \(\mathfrak R_{n,d,\epsilon}=\Theta(n^{-1})\).  Conversely, suppose \(\mathfrak R_{n,d,\epsilon}=O(n^{-1})\).  If \(d>n^2\) along a subsequence, (5) is bounded away from zero, a contradiction.  Otherwise \(d\le n^2\) eventually, and (4) gives \(c_0\sqrt d/n=O(n^{-1})\), hence \(d=O(1)\).  This proves the equivalence and also identifies exactly what the displayed bracket leaves unresolved.
\end{proof}
\par\smallskip\noindent \textit{Justification.} This theorem is the paper's finite-sample minimax bracket: it joins two explicit converse mechanisms to a computable empirical-ratio upper bound and derives the exact proved consistency and parametric-risk consequences. \textit{Gap.} The accepted in-repository linear-ATE theorem \texttt{stat\_discrete\_ate\_minimax\_loggap/thm:sharp-minimax-fixed-interior} matches \(n^{-1}+d^2/[n^2\log^2 n]\) with a heavy/light hybrid. Here the nonlinear maximum creates a tie singularity, lower mechanisms \(\sqrt d/n\) and \(d/[n\log(en)]\), and only \(d/n\) achievability, so neither theorem implies the other. Manski2004 and Stoye2009 instead optimize worst-case welfare regret of a treatment allocation rule after finite samples; that decision criterion neither yields nor is yielded by squared-error estimation of the scalar observational oracle value. Luedtke--van der Laan is local regularity theory, Whitehouse et al. impose gap-density smoothing conditions, and Jiao--Han--Weissman supplies only the normalized-distance lower mechanism after embedding. The crossing at \(d\asymp\log^2(en)\) is not a matched minimax phase boundary. \textit{Consumer.} The theorem tells downstream analysts when consistency is impossible, when the explicit estimator is consistent, and exactly which rate gap remains, without conflating that point-estimation task with finite-sample treatment-rule regret.

\section{Interpretation}

The result is a finite-sample minimax bracket, not a solved minimax equivalence. In the nonsaturated lower bound, \(\sqrt d/n\) and \(d/[n\log(en)]\) are equal at \(d\asymp\log^2(en)\); this is only the crossing of two proved converse mechanisms. The computable empirical-ratio estimator guarantees consistency when \(d=o(n)\), whereas the converse makes \(d=o(n\log n)\) necessary. This is point estimation of a scalar observational oracle value, not the randomized-rule welfare-regret criterion of Manski2004 or Stoye2009: neither task determines the other. The region between those conditions, and an estimator matching either lower mechanism, is unresolved.

\section*{Technical internal limitation (diagnostic only)}

A root-changed matching-achievability attempt was made at the corrected target \(\min\{1,\sqrt d/n+d/[n\log(en)]\}\). Put \(L=\log(en)\), \(K=\max\{1,\lfloor\epsilon^2L/2^{16}\rfloor\}\), and use the frozen three-way split. The decidable pilot sends a cell to the ratio branch exactly when \(\widehat p_x^{(0)}\ge256L/m_0\), both pilot arm counts are at least \(64L\), and \(\lvert\widehat\Delta_x^{(0)}\rvert\ge32\sqrt{L/(m_0\widehat p_x^{(0)})}\); every sparse or near-diagonal cell, including a heavy near-tie cell, goes to the polynomial branch. Decompose the target as one half the sum of \(\Lambda\) and \(\Gamma\). The accepted linear-ATE light/heavy calculations adapt to the two signed linear ratio components of \(\Lambda\), but not to \(\Gamma\). For \(\Gamma\), the attempted branch searches over \(M_K=\binom{K+4}{4}-1\) coefficients \(c_r\) in \(P_x(u)=\sum_{1\leq\lvert r\rvert\leq K}c_ru^r\), with the zero anchor \(P_x(0)=0\), on the compact pilot box \(B_{x,n}\). It accepts the lexicographically first rational certificate only if \(\sup_{u\in B_{x,n}}\lvert P_x(u)-\Gamma(u)\rvert\le2^{12}\epsilon^{-3}b_x/K\) and \(\sum_r\lvert c_r\rvert b_x^{\lvert r\rvert}\le b_x(64/\epsilon)^K\), where \(b_x=2\sum_{a,y}\widehat q^{(0)}_{ay,x}+128L/m_0\); sign splitting and clearing the positive arm sums reduce this finite-dimensional feasibility check to semialgebraic quantifier elimination. Failure of the certificate triggers the empirical-ratio fallback. On acceptance, the fixed-sample lift is \(\sum_r c_r\prod_{a,y}(N^{(2)}_{ay,x})_{r_{ay}}/(m_2)_{\lvert r\rvert}\), with the convention that terms of degree exceeding \(m_2\) are zero. The first unclosed inequality is uniform feasibility of these two simultaneous pointwise-approximation and absolute-coefficient bounds for every overlap-cone box touching the vertex and for every heavy near-tie box. A Jackson-type error bound alone does not control the absolute coefficients, while the accepted reciprocal Chebyshev certificate controls the smooth linear ratio functional but not the composed absolute contrast. Without this inequality, the within-cell factorial second moment, the cross-cell covariance sum, and the pilot-error contribution cannot be bounded at the target rate. Therefore the attempted selector remains uncertified, and the proved upper bound stays \(C_\epsilon d/n\).

\section{Honest scope}

The causal target is the scalar value of the unrestricted oracle policy, identified under consistency, conditional exchangeability, binary treatment and outcome, iid sampling, and fixed overlap on supported cells. The contribution is the explicit finite-sample minimax bracket, the exact lower reductions, and the impossibility witness. It is not learned-policy regret or the finite-sample randomized-allocation minimax-regret criteria of Manski2004 and Stoye2009: a low-regret rule need not estimate the scalar value, and a scalar estimate does not determine a treatment allocation or its regret. It is also not fixed-dimensional honest confidence-interval adaptation, a matched global rate, or a claim that \(d\asymp\log^2(en)\) is an established minimax phase boundary. Vanishing overlap, growing treatment count, matching achievability, the intermediate consistency region, and exact leading constants remain outside the proved result.

\begin{thebibliography}{99}

  \bibitem{ZengBalakrishnanHanKennedy2024Discrete} Zeng, Zhenghao, Sivaraman Balakrishnan, Yanjun Han, and Edward H. Kennedy (2024, revised 2026). "Causal inference with high-dimensional discrete covariates." arXiv:2405.00118v3.

  \bibitem{JiaoHanWeissman2018L1} Jiao, Jiantao, Yanjun Han, and Tsachy Weissman (2018). "Minimax estimation of the \(L_1\) distance." IEEE Transactions on Information Theory 64(10): 6672--6706; accessible five-page author version and arXiv:1705.00807v7.

  \bibitem{LuedtkeVanderLaan2016OptimalValue} Luedtke, Alexander R., and Mark J. van der Laan (2016). "Statistical inference for the mean outcome under a possibly non-unique optimal treatment strategy." Annals of Statistics 44(2): 713--742.

  \bibitem{WhitehouseChenAusternSyrgkanis2025} Whitehouse, Justin, Qizhao Chen, Morgane Austern, and Vasilis Syrgkanis (2025, revised 2026). "Inference on optimal policy values and other irregular functionals via softmax smoothing." arXiv:2507.11780.

  \bibitem{WuYang2016Entropy} Wu, Yihong, and Pengkun Yang (2016). "Minimax rates of entropy estimation on large alphabets via best polynomial approximation." IEEE Transactions on Information Theory 62(6): 3702--3720.

  \bibitem{JiaoVenkatHanWeissman2015Functionals} Jiao, Jiantao, Kartik Venkat, Yanjun Han, and Tsachy Weissman (2015). "Minimax estimation of functionals of discrete distributions." IEEE Transactions on Information Theory 61(5): 2835--2885.

  \bibitem{CaiLow2011L1} Cai, T. Tony, and Mark G. Low (2011). "Testing composite hypotheses, Hermite polynomials and optimal estimation of a nonsmooth functional." Annals of Statistics 39(2): 1012--1041.

\end{thebibliography}

\end{document}